\documentclass[../main.tex]{subfiles}
\begin{document}
%==========================================TIÊU ĐỀ TẬP========================================
\begin{table}[h]
  \begin{adjustwidth}{-1cm}{}
    \begin{tabular}{>{\centering\arraybackslash}p{6cm}|>{\raggedright\arraybackslash}p{12.5cm}}
      \multirow{5}{6cm}
      % Cột bên trái: ÉPISODE và số tập
      {\\[-40pt]
        \flaregothic\fontsize{20pt}{20pt}\selectfont \centering
        {\contour{errorfive}{\textcolor{black}{ÉPISODE}}}\color{black}\\
        \burbank\fontsize{72pt}{72pt}\selectfont
        \includegraphics[height=3.12359cm]{../../../../Image/Book?/number/num38.png}
        \\[18pt]
        % \flaregothic\fontsize{14pt}{14pt}\selectfont
        % \color{epattr}BIENVENUE, \uline{\color{Banpire} Mel} !
        \flaregothic\fontsize{18pt}{18pt}\selectfont{\contour{errorfive}{\textcolor{white}{FIN DE TOME\dots\ !}}}\\
        \fontsize{12pt}{12pt}{\contour{errorfive}{\textcolor{white}{(pour de vrai)}}}
        \color{black}
      }
       &
      % Dòng thứ nhất: tên tập tiếng Nhật
      {\fotskip\fontsize{18pt}{18pt}\selectfont さようなら、\ruby{青}{あお}\ruby{上}{がみ}}                            \\[6pt]
      % Dòng thứ hai: tên tập tiếng Anh
       & {\cafeteria\fontsize{18pt}{18pt}\selectfont See You Again, Aogami!}                               \\[4pt]
      % Dòng thứ ba: tên tập tiếng Việt
       & {\vnmsans\fontsize{18pt}{18pt}\selectfont Hẹn gặp lại, Aogami!}                                   \\[8pt]
      % Dòng thứ tư: nhân vật xuất hiện
       & {\caviar{\fontsize{14pt}{14pt}\selectfont Nhân vật: Tất cả các nhân vật xuất hiện từ đầu series}}
      % \\[6pt]
      % Dòng cuối cùng: Thời gian diễn ra của tập (thường sẽ là ngày phát hành)
      % & {\continuum\fontsize{16pt}{16pt}\selectfont Ngày phát hành: 11/06/2022}\\[8pt]
    \end{tabular}
  \end{adjustwidth}
\end{table}
%===========================================NỘI DUNG==========================================
\color{black}
\vspace{32pt}

\pov{Hikawa Kuon}

\begin{center}
  \uline{(Tiếp tục từ {\flaregothic ÉPISODE \frank 37}.)}
\end{center}

Tất cả chỉ còn lại một đêm nay. Một đêm cuối cùng.

Tôi đứng giữa sân đền Aogami, đôi chân như bị đóng đinh vào mặt đá. Tiếng gió đêm lùa qua hàng đèn lồng khiến chúng khẽ lắc lư, ánh sáng cam nhảy múa trên khuôn mặt đông cứng của mọi người xung quanh tôi. Không ai nói gì. Không ai thở mạnh. Chỉ có tiếng gió và tiếng tàn tro pháo hoa rơi lất phất xuống mặt đất.

Sáu tiếng đồng hồ.

Cái cụm từ đó vang vọng trong đầu tôi như một bản án.

Chúng tôi vừa mới hoàn thành Nghi thức Thanh tẩy cho Naomi, vừa mới nếm được vị ngọt của chiến thắng, vừa mới ôm lấy cơ thể ấm áp của con bé và cảm nhận nhịp tim đập vững vàng lần đầu tiên\dots\ thì lập tức, một bản tin đỏ từ Game đập sầm vào mặt tất cả chúng tôi.

``\eng{Sáu tiếng, Kuon-user. Không phải sáu ngày, càng không phải sáu tháng. Là đúng sáu tiếng. Các cậu không nghe nhầm đâu.}''

Xung quanh tôi, lớp 1-C bắt đầu phản ứng. Không phải theo kiểu hỗn loạn, mà là kiểu chết lặng. Chậm rãi. Từng người một, như cách mà tin dữ ngấm vào da thịt.

Honoka là người đầu tiên lên tiếng, giọng run rẩy: ``Sáu tiếng\dots\ là sáu tiếng thật sao?''

Ayame siết chặt chiếc quạt giấy trong tay, bàn tay trắng bệch dưới ánh trăng. ``Bọn họ\dots\ phải đi trước khi mặt trời mọc ư\dots?''

Shion đứng cách tôi vài bước, mắt vẫn còn ngấn nước vì kiệt sức sau nghi thức. Cô hỏi tôi bằng giọng nghiêm trang, không còn dấu vết của phong thái chuunibyou: ``Kuon\dots\ không có cách nào trì hoãn sao?''

Tôi lắc đầu. ``Game nói rằng năng lượng linh mạch Aogami cộng hưởng với năng lượng ngoại lai trong cơ thể chúng tôi. Cánh cổng sẽ mở lúc bình minh, tồn tại mười lăm phút, rồi biến mất. Không ai biết khi nào nó sẽ xuất hiện trở lại. Chắc chắn\dots\ chúng tôi sẽ không muốn bỏ lỡ đâu.''

Một sự im lặng nặng nề sập xuống.

Rồi, từ đâu đó trong đám đông, tiếng Touka vang lên, nhỏ nhưng rõ ràng: ``Vậy\dots\ chúng ta phải làm gì trong sáu tiếng còn lại?''

Câu hỏi đó như một que diêm châm vào đống củi khô. Mọi người nhìn nhau. Mắt ai cũng đỏ hoe, nhưng không ai khóc nữa. Không còn thời gian để khóc.

Tôi hít vào thật sâu, rồi thở ra chậm rãi.

\ct

Nhưng trước khi nghĩ tới chuyện chia tay, tôi phải giải quyết một vấn đề lớn hơn nhiều.

Naomi.

Con bé đang ngồi trên bậc đá trước cổng torii, đầu gục vào vai Suzuki, mắt sưng húp lên sau khi phải trải qua cơn đau đó. Cơ thể mới của em vẫn còn yếu ớt sau nghi thức, bàn tay nhỏ xíu nắm chặt lấy vạt áo chị gái nuôi như sợ buông ra sẽ mất.

Tôi quỳ xuống trước mặt con bé, đặt tay lên đầu em. Mái tóc mềm mại, ấm nóng, đang sống. Nhịp tim đập đều đặn truyền qua lớp vải mỏng khi tôi chạm vào vai em.

``Naomi này.''

Con bé ngẩng lên. Đôi mắt lục bảo ấy nhìn thẳng vào tôi, trong veo như hai viên ngọc, nhưng chứa đầy nỗi sợ hãi mà một đứa trẻ tám tuổi không nên phải gánh chịu.

Tôi cố giữ cho giọng mình thật bình thản: ``Em nghe anh nói. Cánh cổng sẽ mở lúc sáu giờ sáng. Anh, chị Kaigou, và chị Suzuki\dots\ sẽ phải đi qua cánh cổng đó để trở về thế giới của bọn anh.''

``\dots Em biết rồi.'' Giọng con bé lạc đi.

``Cơ thể của em được tạo ra từ linh khí của Aogami,'' tôi tiếp tục, từng chữ như bị mài bằng đá. ``Game đã phân tích rồi đấy. Nếu em đi qua cánh cổng, xác suất linh hồn của em\dots\ bị phân rã\dots\ là rất cao. Bọn anh không thể mạo hiểm mạng sống mà bọn anh vừa mới giành giật lại cho em được.''

Tôi ngừng lại. Nước mắt con bé lại lăn dài.

``Anh muốn em ở lại đây, Naomi. Ở lại với chị Nanase, với tất cả mọi người. Bọn họ sẽ là gia đình mới của em. Em sẽ được sống, được lớn lên, được--''

``\textbf{Mọi người lại bỏ rơi em!}''

Tiếng hét của Naomi cắt đứt lời tôi, vang vọng khắp sân đền.

Con bé đứng bật dậy, nước mắt tuôn xối xả, hai bàn tay nhỏ xíu nắm chặt thành nắm đấm. Cơ thể bé bỏng run lên bần bật, không phải vì lạnh, mà vì một cơn giận dữ xen lẫn tuyệt vọng.

``Anh nói không phải bỏ rơi\dots\ \textbf{nhưng cuối cùng vẫn là bỏ em lại!}'' con bé gào lên, giọng nghẹn ngào. ``Trước đây em cũng bị như vậy! Lúc em còn là một linh hồn, không ai nghe thấy em, không ai nhìn thấy em! Em gọi mãi, gọi mãi mà không ai trả lời cả! Em tưởng\dots\ em tưởng khi có anh, có chị Kaigou, có chị Suzuki\dots\ em sẽ không bao giờ bị bỏ lại một mình nữa\dots\ Vậy mà giờ\dots anh chị lại\dots''

Lời nói của Naomi không dứt. Con bé sụp xuống đất, ôm mặt khóc nức nở. Nhìn khung cảnh đáng thương đó, tôi chỉ biết đứng chôn chân, không nói được gì.

Không phải vì tôi không có câu trả lời. Mà đúng hơn\dots\ câu trả lời duy nhất tôi có lại là đáp án phũ phàng nhất. Không một ai muốn bị tách rời khỏi những người thân yêu của mình cả, tôi hiểu cảm giác đó. Nhưng\dots\ tôi không thể hành động một cách mù quáng chỉ vì tình cảm cá nhân được. Tôi phải suy nghĩ cho tương lai của con bé, nghĩ tới những gì tốt nhất cho con bé khi chúng tôi đã rời đi. Mong rằng Naomi có thể hiểu được điều đó, dù tôi biết rằng nó sẽ rất khó khăn lúc đầu.

``\dots''

Tôi nhìn em ấy, nhìn đôi vai bé nhỏ đang rung lên vì khóc, và trong đầu tôi, một hình ảnh khác chồng lên.

Không phải ở đây. Không phải ở sân đền Aogami. Mà ở \emph{The Void.} Hồi tôi còn một mình lạc lối tại đó.

Một không gian đen đặc, không ngày đêm, không phương hướng. Và giữa bóng tối tưởng chừng như vô tận ấy, hai cô gái cũng đã từng đứng trước tôi với cùng một ánh mắt: ánh mắt của những con người bị thế giới bỏ rơi.

Hikawa Kaigou, một cô gái hai mươi lăm tuổi, gánh cả cuộc đời trên vai từ năm mười tám, mất cả cha lẫn mẹ, bị họ hàng quay lưng, bị xã hội dồn vào góc tường. Cô đã sống bằng nỗi hận thù và sự cảnh giác, xây lên bức tường đá chắn giữa mình và cả thế giới, chỉ để bảo vệ một cô em gái duy nhất.

Hikawa Suzuki, mười lăm tuổi, học chậm hai năm vì thiếu tiền, ăn mì gói cả tuần, ngủ trên tấm đệm cũ sờn. Cô bé từng nói với tôi: ``Nếu em cũng gục xuống, thì Onee-sama biết bám vào đâu?''

Và khi tôi đề nghị họ đi cùng tôi, rời bỏ thế giới cũ để bắt đầu một cuộc đời mới, Kaigou lúc đầu đã từ chối. Không phải vì không muốn, mà vì \emph{không dám tin}, nhất là từ một người con trai như tôi. Cô ấy đã chịu tổn thương quá nhiều rồi, cô không thể mở lòng để tiếp tục chịu những lời nói mật ngọt đầy sự dối trá thêm một lần nào nữa. Phải đến khi chính cô em gái mà cô bảo vệ suốt đời tuyên bố rằng ``em sẽ đi theo anh ấy, dù có chị hay không,'' Kaigou mới sụp đổ hoàn toàn.

Cô đã khóc. Lần đầu tiên sau bảy năm.

Và cô đã chọn đi cùng tôi.

Tôi nhớ cảm giác lúc đó. Cảm giác khi một con người quyết định trao toàn bộ niềm tin còn sót lại cho một kẻ xa lạ, chỉ vì không còn gì để mất nữa. Đó không phải là sự tin tưởng, đó là sự liều mạng, một canh bạc cuối cùng của những kẻ cùng đường.

Nhưng với Naomi\dots\ mọi chuyện khác hẳn.

Kaigou, Suzuki, và cả Game nữa, tất cả đều bị thế giới cũ vứt bỏ. Thế giới đó không còn gì cho họ. Với hai chị em, đó là căn phòng trọ cháy rụi, khu chung cư mục nát và một xã hội bất công. Với Game, đó là sự ruồng bỏ của người tạo ra mình, bị lãng quên trong những góc tối bộ nhớ với một mã nguồn dở dang và một trái tim đầy lỗi, cùng nỗi đau lạc mất người thương. Khi tôi mời họ đi, tôi đang kéo họ ra khỏi một nơi chết tiệt để đến một nơi tử tế hơn. Đó là một phép tính đơn giản: mất ít, được nhiều.

Naomi ban đầu cũng như vậy. Cô bé bị sự mu muội của người lớn nhẫn tâm cướp đi mạng sống bằng nghi thức hiến tế và không được ai công nhận suốt 1300 năm. Nhưng giờ đây, em không bị thế giới này bỏ rơi. Aogami là nơi em được sinh ra một lần nữa. Linh khí của mảnh đất này chảy trong cơ thể em. Lớp 1-C yêu thương em. Shino, Nanase, Ayame, tất cả đều sẵn sàng trở thành gia đình mới của em.

Nếu tôi đưa em đi, tôi không phải đang cứu em khỏi một thế giới tàn nhẫn. Tôi đang kéo em ra khỏi nơi duy nhất mà em thuộc về, để đặt vào một canh bạc sinh tử mà ngay cả Game cũng không dám đảm bảo kết quả có lợi.

Đó chính là sự khác biệt căn bản. Tôi không thể vác công thức dùng cho hai chị em đó áp vào cô bé được.

Tôi cúi đầu, nhắm mắt lại. Hình ảnh của Kaigou khóc trong The Void và Naomi khóc trước cổng torii chồng lên nhau, nhưng chúng không giống nhau. Chúng giống nhau ở bề ngoài: đều là nỗi đau bị bỏ rơi. Nhưng bản chất thì hoàn toàn trái ngược.

Với Kaigou, tôi mở một cánh cửa.

Với Naomi, tôi có thể đang đóng một cánh cửa.

Và đó là lý do mà tôi không thể dễ dàng gật đầu.

Tôi buông một tiếng thở dài nặng nề, quỳ sát lại và đặt hai tay lên vai em.

``Naomi\dots\ Nếu em đi với anh, chưa chắc em đã có được sự bình yên như thế này đâu. Có thể nó sẽ rất ồn ào, phức tạp, và đầy những mối nguy hiểm,'' tôi nhẹ giọng giải thích, hy vọng rằng em ấy sẽ hiểu ý. ``Nơi này\dots\ Aogami này là nơi tốt nhất cho em, một nơi thanh bình và tuyệt đẹp. Em sẽ được sống một cuộc đời bình thường mà em xứng đáng có.''

Nhưng rồi, ngay sau khi tôi dứt lời, Naomi liền khẽ lắc đầu. Con bé dùng mu bàn tay quệt ngang dòng nước mắt, ngước lên nhìn tôi. Đôi mắt lục bảo từng chứa đựng ngàn năm tĩnh lặng nay lại rực cháy một ánh sáng bướng bỉnh và kiên định đến lạ lùng.

``Thế giới có đẹp và yên bình đến mấy mà không có Onii-chan và Onee-chan, thì với em đó cũng chỉ là một vòng lặp cô đơn khác mà thôi.''

Câu nói đó\dots\ như một mũi tên ghim thẳng vào tim tôi.

Một vòng lặp cô đơn.

Phải rồi. Sống một mình trong một thế giới hoàn mỹ, dẫu có hoa thơm cỏ lạ, dẫu có những người bạn tốt bụng\dots\ nhưng lại thiếu vắng đi những người quan trọng nhất, những người đã trao cho em hơi ấm đầu tiên, thì cảm giác đó chẳng khác nào bị giam lỏng trong một cái lồng son.

Tôi nhớ lại ngày đầu tiên chúng tôi đưa em đi dạo phố, cái hôm sau khi em vừa được tái sinh. Cái ngày mà em đã chạy lăng xăng trên con đường lát đá dưới ánh nắng rực rỡ của mùa hè, trầm trồ trước ly parfait trà xanh lạnh buốt, sợ hãi nép vào lưng Suzuki khi thấy xe cộ rồi lại reo lên thích thú khi thấy đàn cá Koi dưới chân cầu.

Ngày hôm đó, khi hoàng hôn buông xuống nhuộm cam cả công viên, em đã ngả đầu tựa vào vai Suzuki và thì thầm: ``Một gia đình sao\dots\ Em cũng mong có một gia đình ấm áp như thế này.''

Và tôi nhớ cả lời khẳng định kiên quyết của em với Suzuki vào một đêm trong những ngày mà chúng tôi phải chạy đua với thời gian vì ôn thi. Lúc đó, tôi đã nghe được lời hứa mà tôi đã tình cờ nghe được ở ngoài cửa phòng: ``Dù sau này có chuyện gì xảy ra, dù chúng ta có đi tới đâu, em vẫn muốn được sát cánh cùng anh chị đến tận cùng ạ.''

Em đã coi ba kẻ ngoại lai chúng tôi là gia đình. Một gia đình thực sự.

Và nguyên tắc duy nhất của một gia đình\dots\ là không bao giờ bỏ rơi nhau.

Tôi ngước nhìn Kaigou. Cô nàng khoanh tay trước ngực, đôi môi mím chặt đến mức tái nhợt, đôi mắt nhìn chằm chằm xuống nền đá như đang cố gắng đấu tranh giữa lý trí và tình cảm.

``Kuon-kun nói đúng đấy, Nao,'' Kaigou lên tiếng, giọng cô run rẩy nhưng cố gắng giữ sự lạnh lùng lý trí để làm điểm tựa. ``Linh hồn em quá mỏng manh. Cánh cổng không gian không phải là một trò đùa, bọn chị đã trải qua nó ba lần rồi và lần nào cũng như một chuyến tàu tử thần. Chị\dots\ chị không muốn nhìn thấy em tan biến thêm một lần nào nữa. Chị thà để em ở lại đây, biết rằng em vẫn sống, vẫn thở, còn hơn là mất em vĩnh viễn trong một khoảng không vô định.''

Kế bên, Suzuki đã ôm mặt khóc nức nở. Cô bé vốn là người chiều chuộng Naomi nhất, là người đóng vai trò ``người mẹ'' chăm sóc em từng bữa ăn giấc ngủ, giờ đây lại không thể thốt nên lời. Em ấy hiểu lý lẽ của tôi và Kaigou, hiểu sự tàn khốc của việc di chuyển xuyên thứ nguyên, nhưng trái tim cô thì gào thét không muốn rời xa sinh mệnh bé nhỏ này.

``Xác suất để thực thể Naomi sống sót khi đi qua cổng không gian\dots\ là rất thấp,'' Game bất ngờ lên tiếng, giọng nói trầm khàn vang lên từ chiếc đồng hồ, phá vỡ bầu không khí bế tắc. ``Tuy nhiên\dots\ nó không phải là 0\%.''

Tất cả ánh mắt, từ nhóm Hikawa chúng tôi cho đến cả lớp 1-C, đồng loạt đổ dồn vào cổ tay tôi.

``Nghi thức Thanh tẩy vừa rồi không chỉ neo giữ cơ thể vật lý, mà còn cường hóa cấu trúc linh hồn của cô bé bằng năng lượng thuần khiết nhất của hệ thống tâm linh Aogami,'' Game phân tích tiếp, từng thông số kỹ thuật được dịch ra ngữ mà chúng tôi có thể hiểu được. ``Nếu được bảo vệ trong trường năng lượng ngoại lai của ba người các cậu khi đi qua cổng, cộng thêm ý chí sinh tồn mãnh liệt của chính cô bé\dots\ hệ thống tính toán rằng cô bé có một cơ hội mỏng manh để vượt qua ranh giới thứ nguyên mà không bị phân rã. Nhưng tôi nhắc lại: đó là một canh bạc khổng lồ. Quyết định cuối cùng\dots\ không nằm ở sự tính toán của tôi hay sự cho phép của các cậu, mà nằm ở chính thực thể này. \uline{Naomi phải tự mình quyết định xem cô bé có muốn đi hay không.}''

Tôi nhìn xuống Naomi. Con bé không hề tỏ ra sợ hãi trước từ ``tan biến'' hay ``canh bạc''. Bàn tay nhỏ xíu của em nắm lấy vạt áo tôi càng lúc càng chặt hơn, đôi mắt ướt đẫm nước mắt nhưng rực sáng ngọn lửa của sự kiên định.

``Em biết là nguy hiểm,'' Naomi thút thít, ngước cằm lên. ``Nhưng em muốn được ở bên cạnh mọi người. Dù chỉ thêm một giây thôi cũng được. Nếu phải tan biến\dots\ thì tan biến cùng nhau vẫn tốt hơn là sống sót một mình ở đây. Xin anh đấy, Onii-chan\dots\ cho em đi cùng đi.''

Cả sân đền Aogami chìm vào một sự tĩnh lặng thiêng liêng. Những học sinh lớp 1-C chứng kiến toàn bộ câu chuyện từ nãy đến giờ, không một ai lên tiếng phản đối. Bọn họ dường như thấu hiểu được nỗi khao khát có một mái ấm của cô bé mang trong mình lời nguyền một ngàn ba trăm năm này.

Shion bước lên một bước, chống tay lên hông, khẽ mỉm cười; một nụ cười dịu dàng và trưởng thành nhất mà tôi từng thấy từ cô nàng Pháp sư vốn luôn đắm chìm trong ảo tưởng chuunibyou.

``Kuon-san. Phong ấn của tôi và Sakura không phải là đồ bỏ đi đâu,'' cô nháy mắt, gõ nhẹ cây quyền trượng gỗ xuống đất. ``Linh hồn của con bé giờ đã vững vàng hơn những gì cậu tưởng tượng rất nhiều. Nếu con bé đã có dũng khí để đánh cược cả sinh mạng\dots\ thì cậu với tư cách là một `Onii-chan' không có quyền tước đoạt đi dũng khí đó.''

Sakura đứng bên cạnh gật đầu đồng tình, nắm chặt lấy tay Koishin. Saki, Miku, và cả những người khác trong lớp cũng khẽ gật đầu, nở nụ cười khích lệ. Bọn họ, dù rất muốn Naomi ở lại, nhưng họ chọn cách tôn trọng quyết định của trái tim con bé.

Tôi thở ra một hơi dài, cảm giác như trút được một gánh nặng ngàn cân, nhưng đồng thời lại tự gánh lên vai một trách nhiệm còn lớn lao và đáng sợ hơn thế gấp vạn lần. Nếu con bé tan biến trên đường đi, tôi sẽ không bao giờ tha thứ cho bản thân mình.

Nhưng nếu tôi từ chối, tôi sẽ tự tay bóp nghẹt tâm hồn em.

Tôi đưa tay lên, vuốt những lọn tóc đen dính bết vì mồ hôi và nước mắt của Naomi, rồi nhẹ nhàng kéo em vào lòng.

``Được rồi,'' tôi thì thầm, siết vòng tay lại. ``Chúng ta sẽ cùng đi. Sẽ không ai bị bỏ lại phía sau cả.''

Naomi òa khóc nức nở, vòng hai cánh tay nhỏ xíu ôm chầm lấy cổ tôi. Lần này không phải là tiếng khóc của sự hoảng loạn tuyệt vọng, mà là tiếng khóc của sự giải tỏa, của niềm hạnh phúc vỡ òa. Kaigou cũng không kìm nén được nữa, cô lao tới ôm lấy cả tôi và Naomi, những giọt nước mắt nóng hổi làm ướt đẫm vai áo yukata của tôi. Suzuki nhào vào theo, khóc rống lên như một đứa trẻ, tạo thành một vòng tay chặt chẽ của bốn mảnh ghép lạc lõng nhưng đã trọn vẹn tìm thấy nhau.

Nhìn cảnh tượng ấy, không khí nặng nề và ngột ngạt của sự chia ly từ nãy đến giờ dường như đã được xua tan đi một nửa, nhường chỗ cho một sự ấm áp đến nao lòng.

``Thật là\dots\ làm người ta cảm động muốn khóc luôn mà,'' Miku sụt sùi, lấy tay áo yukata chùi đi những giọt nước mắt đang chảy dài trên má, nhưng trên môi lại nở một nụ cười rạng rỡ.

``Bọn họ đúng là một gia đình thực sự,'' Yae khẽ chắp tay trước ngực, ánh mắt đong đầy sự ngưỡng mộ xen lẫn chút u buồn. ``Dù không chung dòng máu, nhưng sợi dây liên kết của họ còn bền chặt hơn bất cứ thứ gì.''

Shino đứng lặng yên bên cạnh, đôi mắt dịu dàng nhìn bốn người chúng tôi đang ôm chặt lấy nhau. Cô khẽ mỉm cười, một nụ cười thấu hiểu sâu sắc từ một người đã từng trải qua nỗi đau mất mát: ``Đúng vậy. Đã là gia đình\dots\ thì dù có đi đến tận cùng thế giới, cũng không thể tách rời được.''

Những tiếng nấc nghẹn ngào dần được thay thế bằng những cái gật đầu và nụ cười ấm áp của các thành viên lớp 1-C. Nỗi buồn chia ly vẫn còn đó, nhưng nó không còn là sự tuyệt vọng tăm tối nữa, mà đã trở thành một sự trân trọng đầy ý nghĩa.

\ct

Khi những cảm xúc đã dần dịu xuống và tiếng nấc chỉ còn là những nhịp thở đều đặn, tôi từ từ đứng dậy, vỗ nhẹ lớp bụi trên tà áo.

Thời gian vẫn đang tàn nhẫn cạn dần. Sáu tiếng bây giờ chỉ còn lại chưa đầy năm tiếng rưỡi.

Tôi đưa mắt nhìn quanh, bắt gặp những ánh mắt buồn sầu nhưng thấu hiểu của ba mươi sáu con người trong lớp 1-C. Bọn họ đã dành trọn cả một mùa hè rực rỡ cho chúng tôi, và giờ đây, chúng tôi sắp phải rời đi, mang theo cả người mà họ muốn bảo vệ nhất.

``Thời gian không còn nhiều,'' tôi lên tiếng, hắng giọng để cố gắng lấy lại tông giọng điệu bình tĩnh thường ngày, dù nó vẫn còn hơi rè. ``Chúng tôi\dots\ muốn quay về `nhà' một chuyến. Căn hộ của ba đứa ở thị trấn này. Một căn nhà mà có lẽ không phải của chúng tôi, nhưng nó là nơi lưu giữ những kỷ niệm cuối cùng của chúng tôi tại Aogami.''

Tôi ngập ngừng một chút, lướt mắt qua từng gương mặt quen thuộc. ``Mọi người\dots\ có muốn đến đó không? Để\dots\ tận hưởng nốt những giờ phút cuối cùng này.''

Nanase chớp mắt. Trong thoáng chốc, sự buồn bã trong đôi mắt của cô nàng Tiểu thư tan đi, nhường chỗ cho một nụ cười rạng rỡ và tinh nghịch nở trên môi.

``Một ý kiến tuyệt vời, Kuon-san,'' Nanase bước lên trước, xoay nhẹ chiếc ô giấy trên tay, tà áo yukata màu thiên thanh khẽ bay trong gió đêm. ``Nếu đã là đêm cuối cùng, sao chúng ta không làm một thứ gì đó thật đáng nhớ nhỉ? Một buổi \textit{slumber party} - tiệc ngủ qua đêm thì sao?''

``Tiệc ngủ sao?'' Miku reo lên, hai mắt sáng rực lên như đèn pha. ``Giống như trong mấy bộ phim thanh xuân ấy hả!?''

``Nghe duyệt đấy!'' Shiina đấm nhẹ hai nắm tay vào nhau, cười nhe răng. ``Đêm nay không ai được ngủ! Chúng ta sẽ quẩy đến tận lúc bình minh luôn!''

``Dù sao thì ngày mai cũng là Chủ nhật mà,'' Touka mỉm cười nhẹ nhàng, vuốt lại nếp áo. ``Thức trắng một đêm với bạn bè cũng là một trải nghiệm không tồi.''

Cả lớp 1-C đồng thanh hưởng ứng. Những tiếng hò reo vang lên khắp sân đền, xua tan đi phần nào không khí u buồn nặng trĩu của sự chia ly sắp tới. Dù tất cả đều biết rõ rằng khi mặt trời mọc vào sáng ngày mai sẽ là khoảnh khắc chia tay đau lòng nhất, nhưng ngay lúc này, bọn họ chọn cách trân trọng và bùng cháy trong từng giây phút hiện tại.

Kaigou lau nước mắt, khẽ huých cùi chỏ vào mạn sườn tôi, mỉm cười mỉa mai: ``\vie{Cái căn hộ bé như lỗ mũi đó mà chứa được hơn ba mươi người thì cậu cũng tài đấy, Kuon-kun.}''

``\vie{Cũng chẳng bé đến mức đó đâu,}'' tôi nhếch mép cười đáp trả. ``\vie{Cùng lắm là hỏng hóc gì thì cứ bảo Hội học sinh lo liệu.}''

Suzuki nắm chặt tay Naomi, cả hai chị em cùng bật cười rạng rỡ, những giọt nước mắt vẫn còn đọng trên mi nhưng nụ cười đã nở rộ như hoa mùa hè.

Đêm nay, tại thị trấn Aogami này, một buổi tiệc ngủ đặc biệt nhất trong lịch sử lớp 1-C chuẩn bị bắt đầu.

\ct

Nói là tiệc ngủ, nhưng tôi biết thừa chẳng có ai ở đây thực sự muốn nhắm mắt vào thời khắc này cả.

Khi chúng tôi rời đền Aogami để trở về nhà - nơi mà chúng tôi đã dành hơn bốn tháng sinh sống và gặp mặt - cả một đoàn người rồng rắn kéo theo sau. Những học sinh lớp 1-C vẫn mặc nguyên những bộ yukata lộng lẫy từ đêm lễ hội, lấp lánh dưới ánh đèn đường vàng vọt. Không chỉ có học sinh, một số người dân thị trấn và gia nhân nhà Furukawa cũng đi theo, mang theo đủ thứ từ chăn màn, gối nệm, cho đến những hộp thức ăn còn thừa từ buổi tiệc tàn, và cả mấy thùng đồ uống lạnh buốt.

``Cẩn thận bước chân nhé mọi người!'' Ayame lớn tiếng chỉ đạo ngay từ ngoài cổng lớn, chiếc quạt giấy trên tay được dùng như một cây gậy chỉ huy của nhạc trưởng. ``Sắp xếp giày dép gọn gàng ở huyền quan! Ai không có chỗ ngồi bên trong thì trải nệm tràn ra hiên nhà và ban công nhé! Nhớ chừa lối đi cho mọi người đi vệ sinh hoặc lấy đồ ăn!''

Và thế là, chỉ trong vòng chưa đầy mười lăm phút, căn hộ tĩnh lặng của tôi biến thành một trại tị nạn-- à không, một câu lạc bộ đêm thu nhỏ mang phong cách thanh xuân vườn trường.

Những chiếc đệm futon được trải la liệt khắp sàn nhà, lấn cả ra ngoài hành lang gỗ hẹp. Những ngọn đèn huỳnh quang trắng lạnh lẽo trong phòng được tắt đi, nhường chỗ cho ánh sáng vàng cam ấm áp từ vài chiếc đèn ngủ nhỏ và những chiếc đèn lồng giấy mượn từ đền. Không có nhạc xập xình ồn ào, chỉ có tiếng gió đêm mùa hè rì rào thổi qua tán cây ngoài cửa sổ, tiếng dế kêu rả rích ngoài hiên, và tiếng trò chuyện rầm rì của hơn ba mươi cái miệng cùng một lúc.

``Tuyệt thật đấy!'' Honoka reo lên phấn khích khi cô nàng thành công giành được một góc nệm êm ái cạnh chiếc bàn trà thấp, giơ cao lon nước trái cây. ``Một buổi tiệ ngủ độc nhất vô nhị luôn! Thề là sau này có tốt nghiệp đại học, lấy chồng sinh con đi chăng nữa, tớ cũng không bao giờ quên được đêm nay đâu!''

``Đừng có nói gở lúc này chứ, cái đồ ngốc này,'' Touka ngồi cạnh đó, khẽ gõ nhẹ vào đầu cô bạn bờm Cáo sa mạc, nhưng đôi mắt vẫn đong đầy ý cười.

Yuko-sensei cũng có mặt. Cô giáo chủ nhiệm của chúng tôi không chen chúc vào đám đông ồn ào bên trong, mà lặng lẽ chọn một góc tĩnh lặng ngoài hiên nhà, tựa lưng vào cột gỗ. Cô cởi bỏ đôi guốc cao gót, duỗi chân thoải mái, cầm trên tay một ly trà xanh còn bốc khói nhẹ. Cô chỉ mỉm cười trìu mến nhìn đám học trò của mình đang phá vỡ mọi quy tắc về giờ giới nghiêm, không buông một lời la mắng. Có lẽ, cô cũng hiểu rằng đêm nay là một ngoại lệ không thể lặp lại.

Tôi lách người qua đám đông, cẩn thận bước qua chân của Aku và Sakura đang ngồi dựa vào nhau, tay bê một khay lớn đựng những cốc trà lúa mạch lạnh mà Kaigou vừa pha xong dưới bếp.

``Trà đây,'' tôi đặt khay xuống giữa bàn, nơi nhóm của Suzuki, Suzune, Hanabi và Yuka đang tụ tập vây quanh những bịch snack mở tung. ``Uống đi cho đỡ khát. Đừng có ăn bánh kẹo nhiều quá kẻo sâu răng đấy.''

``Cảm ơn Onii-sama!'' Suzuki vui vẻ nhận lấy một cốc, rồi quay sang thổi nhẹ và đút cho Naomi đang cuộn tròn như một chú sâu nhỏ trong lớp chăn mỏng bên cạnh. Con bé há miệng uống ngoan ngoãn, đôi mắt lục bảo tò mò nhìn mọi thứ xung quanh.

Nhìn quanh căn phòng chật chội nhưng tràn ngập hơi ấm, tôi bất giác mỉm cười. Bầu không khí tĩnh lặng và bi thương ở sân đền Aogami lúc nãy dường như đã bị ép lùi lại, tạm thời nhường sân khấu cho sự náo nhiệt của tuổi thanh xuân. Có một điều gì đó thật sự kỳ diệu ở những bữa tiệc đêm khuya thế này. Nó giống như một khoảng không gian bị cắt rời khỏi dòng chảy thời gian thực, nơi những lo âu về ngày mai bị lãng quên, nơi mọi áp lực bị trút bỏ, chỉ còn lại sự ấm áp của hiện tại.

Mọi người chia thành từng nhóm nhỏ, rải rác khắp các góc phòng. Những tiếng cười nói, những lời trêu chọc xen lẫn những tiếng thở dài hoài niệm. Tôi lùi lại, tựa lưng vào bức tường cạnh cửa sổ, để mặc cho sự mệt mỏi của một ngày dài len lỏi vào cơ thể, thả lỏng các bó cơ đang căng cứng.

Đương nhiên, đã gọi là tiệc ngủ thì không thể thiếu những ``thủ tục'' truyền thống của lứa tuổi học trò. Dù vừa mới trải qua một mùa hè đánh vật với những lời nguyền và oán linh thực sự, nhưng dường như điều đó vẫn không làm giảm đi sự tò mò của họ đối với những thế lực siêu nhiên. Thủ tục đầu tiên được khởi xướng, không gì khác ngoài: kể chuyện ma.

Không biết từ lúc nào, một vòng tròn đã được thiết lập ngay giữa căn phòng sinh hoạt chung. Bốn ngọn đèn lồng mượn từ đền Aogami được đặt ở bốn góc, leo lét tỏa ra ánh sáng cam mờ ảo. Đèn pin điện thoại được bật lên, hắt ánh sáng ngược từ dưới cằm lên khuôn mặt đang cố tỏ ra nguy hiểm và bí ẩn của Honoka. Xung quanh cô nàng, Saki, Hanabi, Yuka và vài nữ sinh khác đang ôm chặt lấy những chiếc gối, cơ thể hơi run rẩy nhưng đôi mắt lại mở to đầy vẻ mong chờ, giống hệt như những người sợ xem phim kinh dị nhưng vẫn lén nhìn qua kẽ tay.

``Tại sao\dots\ chúng ta lại phải chơi trò này lúc một rưỡi sáng cơ chứ?'' Saki rên rỉ, rúc nửa khuôn mặt vào chiếc gối ôm hình con gấu trúc mà cô nàng đã cất công mang từ tận nhà đến. ``Trường mình thiếu gì chuyện rùng rợn thật đâu mà phải kể thêm\dots''

``Vì đây là luật bất thành văn của mọi buổi tiệc ngủ! Nếu không có chuyện ma thì không gọi là tiệc ngủ!'' cô nàng bờm Cáo sa mạc tuyên bố chắc nịch, rồi cố tình hạ thấp tông giọng xuống mức thầm thì rùng rợn, đưa chiếc đèn pin lắc lư qua lại. ``Nghe cho kỹ đây\dots\ câu chuyện này xảy ra ở chính ngôi trường của chúng ta\dots\ ngay tại cầu thang dãy nhà cũ, nơi mà ánh sáng mặt trời không bao giờ chạm tới\dots''

Honoka bắt đầu kể về một nữ sinh từng bị bỏ lại sau giờ học vào một chiều đông lạnh giá. Theo lời kể, cô gái ấy nghe thấy tiếng bước chân vang vọng trên hành lang không người, từ từ tiến lại gần, và rồi\dots\ những tiếng gõ cửa cộc cộc vang lên từ bên trong một phòng học đã bị khóa kín bằng ổ khóa gỉ sét.

Lối kể chuyện của cô nàng thực sự rất cuốn hút, nhịp độ được điều chỉnh lúc trầm lúc bổng, cộng thêm tiếng dế kêu rả rích ngoài hiên nhà và tiếng gió rít qua những khe hở của căn hộ gỗ cũ kỹ vô tình tạo thành một bản nhạc nền hoàn hảo. Khi cô nàng đột ngột vỗ mạnh một cái ``bốp'' xuống chiếu tatami, Saki và Inari gần như thét lên và ôm chầm lấy nhau, nhắm tịt mắt lại.

``Thôi nào hai người,'' Touka ngồi khoanh tay ở vòng ngoài, nhướng một bên lông mày, thẳng tay dội một gáo nước lạnh buốt vào bầu không khí đang căng thẳng tột độ. ``Dưới góc nhìn khoa học, tiếng gõ cửa từ bên trong một căn phòng kín thường là do sự thay đổi áp suất không khí khi nhiệt độ giảm đột ngột vào ban đêm, làm co ngót khung gỗ. Còn tiếng bước chân lộc cộc? Chắc chắn là do kết cấu sàn gỗ cũ bị mọt, tự phát ra tiếng kêu do ứng suất dư thôi. Ma cỏ gì ở đây.''

Cả đám nữ sinh đang chìm đắm trong sự sợ hãi bỗng chốc tụt hết cảm xúc, mặt đần ra như vừa bị giành mất đồ chơi. Miku bĩu môi, phụng phịu tắt chiếc đèn pin: ``Touka-san à, cậu làm hỏng hết không khí rồi! Chuyện ma thì đừng có lôi vật lý, cơ học với kiến trúc vào chứ! Cậu có phải là thợ mộc đâu!''

``Mình chỉ đang cung cấp một góc nhìn duy vật biện chứng để mọi người đỡ sợ thôi mà,'' cô nàng Thần đồng nhún vai, mặt tỉnh bơ không chút hối lỗi, tự tin với lập luận của mình.

``Đôi khi thành thật quá cũng không phải điều tốt lành gì đâu, Touka-san à\dots'' tôi ngồi đối diện Inari, miệng cười chát với lời biện hộ của cô ấy.

Ngồi bên cạnh Touka, Naomi, người từ nãy đến giờ vẫn im lặng lắng nghe, chớp chớp đôi mắt lục bảo trong veo. Con bé nghiêng đầu ngây thơ hỏi, giọng nói tĩnh lặng đến lạ thường: ``Nhưng mà\dots\ cái chị nữ sinh ở hành lang đó\dots\ sao không đi xuyên qua tường để ra ngoài ạ? Hồi trước lúc bị kẹt ở rừng, em toàn đi xuyên qua mấy gốc cây cổ thụ thôi. Có mấy chú lính cụt đầu thời chiến quốc còn hay chui lên từ dưới lòng đất nữa cơ. Bọn họ bảo đi như vậy mới nhanh.''

Sự ngây thơ chân thật của một cựu oán linh ngàn năm tuổi lập tức tạo ra một sự im lặng chết chóc trong phòng. Cả đám nữ sinh nhìn chằm chằm vào Naomi, da gà da vịt nổi rần rần trên tay. So với câu chuyện bịa đặt đầy lỗ hổng logic của Honoka, lời nhận xét tỉnh bơ dựa trên ``kinh nghiệm tâm linh thực tế thời phong kiến'' của Naomi mang một sức sát thương khủng khiếp hơn gấp bội phần.

``\textbf{A-A-Á! Kh-Không kể chuyện ma nữa! Chuyển trò! Chuyển trò ngay lập tức!}'' Saki hét lên thảm thiết, vội vã bò nhoài ra và bật tung tất cả các công tắc đèn trong phòng lên, kiên quyết từ chối lắng nghe thêm bất cứ chi tiết nào về ``kinh nghiệm thực tiễn'' của cô bé. Chúng tôi cười ồ lên trước vẻ hoảng hốt của cô nàng bán hoa, xóa tan đi chút tàn dư lạnh lẽo của câu chuyện ma.

\ct

Trò kể chuyện ma bị dẹp bỏ không thương tiếc, và ngay lập tức, một trò chơi kinh điển khác được thế chỗ: ``Sự thật hay Thử thách''.

Một vỏ chai nước khoáng rỗng được đặt ở giữa chiếc bàn tròn. Luật chơi vô cùng đơn giản và tàn nhẫn: người xoay chai sẽ được quyền đặt câu hỏi hoặc đưa ra thử thách cho người mà đầu chai chỉ vào khi nó dừng lại. Ai từ chối sẽ phải uống một cốc nước ép mướp đắng pha chanh cực chua mà Ayame đã cất công chuẩn bị riêng cho ``những kẻ hèn nhát''.

``Được rồi, tôi bắt đầu trước nhé!'' Shiina xoa hai tay vào nhau đầy hào hứng, rồi dùng sức vung tay xoay mạnh chiếc chai nhựa. Chiếc chai quay tít trên mặt bàn gỗ, tiếng ma sát lạch cạch vang lên thu hút sự chú ý của tất cả mọi người. Sau vài vòng quay chậm dần, đầu chai từ từ chỉ thẳng vào\dots\ Shion.

``Hah!'' Shiina đắc thắng reo lên. ``Nào, Shion. Sự thật hay Thử thách?''

Cô nàng Pháp sư ngồi thẳng lưng, hất lọn tóc màu bạch kim ra sau vai, đôi mắt ánh lên vẻ tự tin thường ngày của một ``Pháp sư bóng tối''. Cô nhếch mép, cất giọng chuunibyou như thường lệ: ``Một người nắm giữ tri thức của vũ trụ thì không có gì phải che giấu. Ta chọn Sự thật.''

Cô bạn bờm Sư tử cười gian xảo, chồm người tới trước: ``Khảo nghiệm nhẹ nhàng thôi. Ở nhà, cậu có hay tự mặc mấy bộ đồ cosplay hắc ám rồi đứng trước gương niệm chú một mình không?''

Sự tự tin trên mặt Shion lập tức vỡ vụn. Đôi má cô nàng đỏ bừng lên như quả cà chua chín. Sakura ngồi cạnh bưng miệng cười khúc khích, trong khi Aku thì giả vờ quay đi chỗ khác.

``Đ-Đó không phải là cosplay!'' Shion ấp úng, cố gắng giữ vớt vát chút thể diện. ``Đó là thánh phục nghi lễ\dots\ và việc thực hành chú ngữ trước gương là để kiểm tra dòng chảy ma lực\dots''

``Nói tóm lại là có đúng không?'' Suzune bật cười, vỗ tay thích thú. Cả phòng lại được một phen cười nghiêng ngả trước sự lúng túng đáng yêu của cô nàng chuunibyou.

Vòng quay tiếp tục. Chiếc chai lại được xoay, và lần này, nó dừng lại trước mặt Nanase. Người xoay là Touka.

``Nanase-san, cậu chọn gì nào?'' cô bạn Thần đồng mỉm cười nhẹ nhàng, nhưng ánh mắt lại ẩn chứa sự ranh mãnh không ngờ.

``Thử thách đi. Một quý tộc không bao giờ chùn bước trước gian nan,'' Nanase nhấp một ngụm trà, phong thái vô cùng điềm tĩnh.

``Vậy\dots\ cậu hãy hát một đoạn bài hát thiếu nhi bất kỳ bằng giọng trẻ con đi.''

Nanase suýt nữa thì phun ngụm trà ra ngoài. Vẻ điềm tĩnh sụp đổ hoàn toàn. ``C-Cậu đùa tôi à? Tôi? Hát nhạc thiếu nhi? Lại còn bằng giọng trẻ con?''

``Đúng vậy,'' Touka gật đầu chắc nịch, chỉ tay về phía cốc nước ép khổ qua xanh rờn sủi bọt trên bàn. ``Hoặc cậu có thể uống cái này.''

Nanase tái mặt nhìn cốc nước, rồi lại nhìn ánh mắt đầy mong chờ của cả lớp. Cuối cùng, cô Tiểu thư danh giá đành nhắm nghiền mắt lại, hít một hơi thật sâu, và cất giọng eo éo hát bài ``Con ếch ộp'' trong tiếng vỗ tay rào rào và những tiếng huýt sáo trêu chọc của mọi người. Sakura ngồi bên cạnh thậm chí còn lấy điện thoại ra quay video lại làm bằng chứng.

Trò chơi tiếp tục quay vòng với những câu hỏi hóc búa và những thử thách dở khóc dở cười. Kanade phải múa bụng trong mười giây. Aku thì bị bắt phải thú nhận lần gần nhất cô trốn học là để đi chơi game thùng. Sự náo nhiệt lên đến đỉnh điểm.

Và rồi, điều gì đến cũng phải đến. Chiếc chai định mệnh dừng lại, và cái đầu chai chỉ thẳng tắp vào ngực tôi.

``Ồ\dots'' Cả phòng đồng thanh ồ lên. Mọi ánh mắt ngay lập tức dồn về phía tôi như những bóng đèn pha rọi vào tội phạm. Người vừa xoay chai là Ayame. Cô nàng Hội trưởng nở một nụ cười nửa miệng đầy mưu mô.

``Đến lượt tổng quản lý rồi,'' Ayame khoanh tay lại. ``Sự thật hay Thử thách đây, Kuon-san?''

Tôi gãi đầu, cảm thấy một chút rắc rối sắp ập xuống. Nếu chọn Thử thách, tôi dư sức làm mọi trò lố lăng mà mọi người nghĩ ra, nhưng với bản tính của Ayame, chắc chắn cô nàng sẽ không bỏ qua cơ hội moi móc thông tin.

``Sự thật đi,'' tôi nhún vai. Dù sao thì tôi cũng không có nhiều bí mật có thể nói ra được, mà những bí mật lớn nhất thì họ cũng đã biết hết rồi.

Ayame híp mắt lại, suy nghĩ vài giây rồi hỏi một câu vô cùng trực diện: ``Kuon-san\dots\ cậu luôn là người che chở và lo lắng cho tất cả chúng tôi, lo cho cả Kaigou-san và Suzuki-chan. Vậy, có bao giờ anh cảm thấy mệt mỏi, và muốn có một ai đó để anh có thể dựa vào không?''

Câu hỏi của Ayame khiến căn phòng đột nhiên tĩnh lặng lại. Không còn những tiếng cười đùa, thay vào đó là sự chú ý cao độ. Ngay cả Kaigou và Suzuki cũng quay sang nhìn tôi, chờ đợi một câu trả lời. Bọn họ biết tôi là Hikawa Kuon của thế giới này, một người đàn ông trưởng thành mang trong mình nhiều kinh nghiệm từ thế giới gốc, luôn khoác lên mình lớp áo giáp gai góc để bảo vệ những người xung quanh.

Tôi cụp mắt xuống, xoay nhẹ lon trà lúa mạch trong tay.

``Có chứ,'' tôi chậm rãi đáp. ``Ai cũng có lúc thấy mệt mỏi. Tôi không phải là siêu nhân, cũng không phải là thần thánh. Nhưng mà\dots'' Tôi ngẩng lên, lướt mắt qua từng gương mặt đang chăm chú nhìn mình. ``Khi nhìn thấy mọi người vẫn an toàn, vẫn có thể cười nói như thế này\dots\ thì cái mong muốn được dựa dẫm đó dường như không còn quan trọng nữa. Nụ cười của mọi người, chính là điểm tựa vững chắc nhất của tôi rồi.''

Một câu trả lời không hề có kỹ xảo, hoàn toàn là những gì chân thành nhất từ tận đáy lòng.

``G-Gì chứ\dots! Cậu sến súa quá rồi đấy!'' Koishin đỏ mặt, quay mặt đi chỗ khác.

``Kuon-san ngầu quá\dots'' Wata lẩm bẩm, mắt rưng rưng.

Câu trả lời của tôi dường như đã làm thỏa mãn trí tò mò của họ, đồng thời cũng khiến bầu không khí lắng xuống một nhịp yêu thương. Trò chơi sau đó lại tiếp tục với người tiếp theo, kéo theo những trận cười mới.

\ct

Đồng hồ điểm hai giờ ba mươi phút sáng.

Sau trò Sự thật hay Thử thách, một vài người đã bắt đầu thấm mệt. Vài tiếng ngáp dài xuất hiện rải rác. Bầu không khí có vẻ đang chùng xuống, báo hiệu cho việc bữa tiệc ngủ sắp chuyển sang giai đoạn\dots\ ngủ thật.

Thế nhưng, có một quy luật bất biến của tuổi trẻ: khi bạn nghĩ mọi thứ đã kết thúc, thì đó mới chỉ là lúc sự hỗn loạn thực sự bắt đầu.

*BỤP!*

Một vật thể mềm mại bay qua không trung và đáp chuẩn xác vào mặt Kaigou khi cô nàng đang lúi húi gọt táo ở góc phòng. Đó là một chiếc gối ôm hình củ cà rốt.

Kaigou từ từ ngẩng mặt lên, mái tóc đỏ rũ rượi, sát khí bừng bừng tỏa ra xung quanh. Cô nhặt chiếc gối lên, nheo mắt nhìn về phía thủ phạm - Shiina, người đang che miệng cười khúc khích vì vừa trượt tay.

``Cậu vừa ném cái gì đấy?'' Kaigou gằn giọng, đôi mắt sắc như dao lam.

``A-Ahaha\dots\ trượt tay! Tôi thề là tôi định ném cho Saki cơ!'' cô bạn bờm Sư tử vội vàng xua tay biện minh.

``Vậy à?'' Kaigou nhếch mép, cân nhắc chiếc gối trong tay. ``Để tôi cho cậu biết thế nào là định luật ba Newton nhé. Lực tác dụng bằng phản lực!''

Và *VÚT!* Chiếc gối bay với tốc độ bàn thờ, đập trúng mặt Shiina khiến cô nàng ngã ngửa ra sau.

``Chơi vậy ai chơi! Yae, hỗ trợ đi!'' Shiina gầm lên, vớ lấy bất cứ chiếc gối nào gần tay và ném trả lại.

Trận đại chiến ném gối chính thức bùng nổ. Không có sự báo trước, không có luật lệ, và chắc chắn là không có sự nhân nhượng.

Những chiếc gối ôm, gối nằm, thậm chí cả thú nhồi bông bắt đầu bay lượn loạn xạ khắp căn hộ như một trận mưa đạn mềm. Phe thể thao với Shiina, Yae, Kaigou lập tức thành lập liên minh tấn công diện rộng. Nhóm Hội học sinh do Ayame chỉ huy phải dựng chăn lên làm chiến hào phòng thủ. Những chiếc gối đập vào nhau bôm bốp. Tiếng la hét, tiếng cười rú lên vỡ òa không kiểm soát.

Và tất nhiên, một người đang đứng ở giữa phòng như tôi không thể nào tránh khỏi việc trở thành nạn nhân của đạn lạc. *BỤP!* Một chiếc gối vuông đập thẳng vào gáy tôi.

Tôi lảo đảo, xoa xoa gáy, quay lại nhìn. Suzuki đang đứng đó, tay cầm một chiếc gối khác, mặt tái mét vội vã chắp tay xin lỗi: ``\vie{Em xin lỗi Onii-sama! Em định ném trúng Koishin-chan cơ!}''

Tôi đứng yên tại chỗ, nhìn khung cảnh hỗn loạn xung quanh. Mọi người đang chạy nhảy, nấp sau những tấm nệm, ném gối vào nhau với sự phấn khích tột độ. Lông vũ từ một chiếc gối bị bục bay lả tả trong không trung như tuyết rơi giữa mùa hè.

Trong khoảnh khắc ấy, ký ức từ ở thế giới thực bỗng ùa về trong tâm trí tôi.

Đó là ngày sinh nhật lần thứ 25 của tôi ở thế giới cũ. Nếu tôi nhớ không nhầm, đó là khoảnh khắc mà các sự kiện chính đã hạ màn và tôi chuẩn bị đi ngủ cho một ngày Chủ nhật mệt mỏi. Nhưng rồi, các cô gái \hll - những đồng nghiệp, những người bạn thân thiết của tôi - đã quyết định tổ chức một bữa tiệc bất ngờ. Và trong đêm đó, các cô gái cũng bày trò ném gối ồn ào y hệt như thế này.

Khi đó, tôi của tuổi 25, gần như thoi thóp bởi núi công việc (thứ mà tôi tự nguyện ôm đồm), cùng với bài vở và đồ án. Tôi chỉ ôm gối, cuộn tròn trong chăn và gắt gỏng yêu cầu họ im lặng để tôi được ngủ. Tôi đã từ chối tham gia vào niềm vui đơn thuần ấy, và rồi bị Watame phá rối, từ đó kéo thành một trận đấu ném gối ``kinh hoàng'' nhất lúc đó, khi cũng chính tôi đã hạ gục Sora - người mà không một ai dám ho he động tới.

Nhưng bây giờ\dots\ tôi không là một tổng quản 25 tuổi cằn cỗi ấy nữa. Tôi là Hikawa Kuon, và tôi đang đứng giữa lớp 1-C, giữa những mảnh ghép quan trọng nhất của cuộc đời tôi ở vũ trụ này. Lần này, thay vì tôi từ chối tham gia vào trò chơi, tôi sẽ là người khuấy đảo trận chiến này lên cao hơn.

*BỐP!* Lại một chiếc gối nữa bay trúng mặt tôi. Lần này là từ Kaigou. Cô nàng cười lớn: ``\vie{Đứng đực ra đấy làm bia đỡ đạn à?}''

Tôi từ từ hạ tay xuống, nhặt chiếc gối dưới chân lên. Một nụ cười nở rộng trên môi, kéo theo những nếp nhăn của sự nhẹ nhõm.

Trái với hồi đó, lần này, tôi không muốn trốn chạy vào giấc ngủ nữa. Tôi muốn thức. Tôi muốn tận hưởng tuổi thanh xuân này, sự ồn ào này, sự ngu ngốc đáng yêu này đến từng giây từng phút cuối cùng.

``\vie{Chơi thế là ăn gian nhé!}'' Tôi gầm lên giả vờ tức giận, hai tay ôm lấy chiếc gối to nhất trên giường và lao vào giữa chiến trường. ``\vie{Tránh đường cho Hikawa Kuon này!}''

Và thế là, tôi tự tay ném đi sự trưởng thành và lớp vỏ bọc nghiêm túc của mình, hòa vào trận chiến với sự nhiệt huyết không kém gì đám học sinh mười lăm tuổi. Tôi ném gối trúng Yae, bị Miku ném trả thù, rồi lại cùng Suzuki lập nhóm phục kích Kaigou. Cả căn hộ rung chuyển vì những bước chân và tiếng hò hét vang dội.

Đó là một đêm không ngủ. Một đêm rực rỡ nhất, ồn ào nhất, và cũng là đêm cuối cùng trọn vẹn nhất mà chúng tôi từng có với nhau.

\ct

Đồng hồ điểm ba giờ sáng.

Sau trận chiến gối ôm kinh thiên động địa, căn hộ nhỏ của chúng tôi trông chẳng khác nào một bãi chiến trường vừa trải qua một cơn bão cấp mười. Lông vũ trắng bay lả tả trên không trung, bám vào tóc và quần áo của từng người. Thế nhưng, trái với sự mệt mỏi thể xác, tinh thần của mọi người dường như đã được gột rửa. Sự nặng nề của cuộc chia ly sắp tới tạm thời lùi lại, nhường chỗ cho những lời tâm tình chân thật nhất.

Tôi biết rõ rằng sự ồn ào và náo nhiệt này chỉ là một lớp vỏ bọc. Dưới những nụ cười rạng rỡ và những câu nói đùa giỡn kia, mỗi người trong số họ đều đang nỗ lực hết sức để kìm nén nỗi buồn chia ly chực trào nơi khóe mắt. Bọn họ gọi đây là một bữa tiệc ngủ, nhưng mục đích thực sự là để kéo dài tối đa thời gian được ở bên cạnh nhau. Họ thức để khắc ghi từng khuôn mặt, từng cử chỉ, từng giọng nói của chúng tôi vào tâm trí trước khi bình minh lạnh lùng mang tất cả đi mất.

``\vie{Nghĩ gì mà trầm ngâm thế, Kuon-kun?}''

Giọng nói quen thuộc mang theo chút mỉa mai nhẹ nhàng vang lên bên cạnh. Tôi quay sang, thấy Kaigou đang cầm một lon soda, tiến lại và tựa lưng vào tường giống hệt tư thế của tôi. Cô nàng đã búi lại mái tóc đỏ rực của mình cho gọn gàng bằng một chiếc kẹp nhựa mượn của ai đó, nhưng đôi mắt vẫn còn lưu lại những vệt đỏ hoe từ trận khóc nức nở ban nãy.

``\vie{Nghĩ xem sáng mai làm sao lẩn đi để không phải dọn dẹp cái bãi chiến trường này,}'' tôi nói dối một cách trơn tru, khẽ nhếch mép.

Cô nàng bật cười, tiếng cười nhẹ nhàng như tan vào tiếng dế đêm, không còn vẻ gai góc thường ngày. ``\vie{Cứ vứt đấy. Mình rời đi qua cổng không gian cơ mà, ai bắt đền được.}'' Cô ngửa cổ uống một ngụm, rồi hướng ánh mắt dịu dàng hiếm thấy về phía Naomi đang được các chị gái trong lớp xúm xít tết lại tóc. ``\vie{Nhưng mà, cậu thấy đấy. Con bé hợp với việc ở bên chúng ta hơn. Dù sao thì\dots\ Nao cũng đã coi chúng ta như một gia đình mà.}''

``\vie{Ừ. Một gia đình thực sự,}'' tôi lặp lại, nhắm mắt lại trong giây lát để cảm nhận sự nhẹ nhõm đang lan tỏa trong lồng ngực. Gánh nặng về sự sống chết dường như đã được tình cảm này san sẻ bớt.

``Kuon-san! Đứng đó làm gì, lại đây đi!''

Tiếng gọi của Wata kéo tôi ra khỏi dòng suy nghĩ. Cậu lớp phó đang vẫy tay rối rít gọi tôi từ phía nhóm của Ayame và Inari đang ngồi quây quần bên góc phòng.

``Đến ngay đây,'' tôi đáp lời, đẩy người khỏi bức tường.

Đêm vẫn còn dài, và mỗi người ở đây đều có những lời muốn nói.

Tôi phủi vài cọng lông vũ trên vai áo yukata, bước về phía góc phòng nơi nhóm của Hội học sinh đang ngồi tĩnh lại sau trận đấu. Ayame đang cẩn thận gắp từng mảnh vụn bánh quy rơi vãi trên chiếu tatami, trong khi Wata thì đang lấy tay xoa xoa vầng trán vừa bị trúng nguyên một cú ném trời giáng. Inari ngồi ngay ngắn bên cạnh, hai tay đan vào nhau đặt trên đùi, ánh mắt trong veo dõi theo từng cử chỉ của tôi.

``Trận chiến ác liệt thật đấy, Kuon-san,'' Ayame ngẩng lên, mỉm cười đưa cho tôi một ly nước lọc. ``Không ngờ cậu cũng có mặt trẻ con như vậy. Tôi cứ tưởng cậu sẽ đứng ngoài khoanh tay thuyết giáo bọn tôi cơ.''

``Thỉnh thoảng cũng phải khởi động gân cốt chứ,'' tôi nhận lấy ly nước, ngồi vắt chéo chân xuống đối diện với họ. ``Dù sao thì tôi cũng là con người, không phải bức tượng đá vô cảm.''

Wata gật gù, ánh mắt nhìn tôi lấp lánh sự ngưỡng mộ không hề che giấu. ``Cậu ném chuẩn xác thật đấy! Quả gối ném vòng cung trượt qua đầu mình đập trúng Yae-san cứ như có mắt vậy! Quả không hổ danh chàng trai độc nhất của lớp 1-C à!''

``Cậu khen tôi hay đang mỉa mai tôi đấy, Wata-san?'' tôi bật cười, lắc đầu. Nhìn ba gương mặt đại diện cho ban cán sự lớp 1-C, giọng tôi bất giác chùng xuống, trở nên nghiêm túc hơn. `Ngày mai chúng tôi đi rồi. Trách nhiệm lèo lái cái lớp cá biệt này sẽ đè nặng lên vai ba người đấy. Mấy người có làm được không?''

Hội trưởng ngừng tay, nụ cười tinh nghịch trên môi cô nàng Hội trưởng dần được thay thế bằng một vẻ chín chắn, điềm tĩnh lạ thường. Cô đưa tay vuốt lại nếp áo, nhìn thẳng vào mắt tôi: ``Kuon-san, cậu coi thường bọn tôi quá đấy. Tôi là Hội trưởng hội học sinh của cả trường Aogami này mà. Mùa hè này, chính cậu đã giúp bọn tôi nhận ra giá trị của sự đoàn kết, giúp bọn tôi phá vỡ những rào cản vô hình giữa những cái tôi quá lớn. Quãng thời gian qua\dots\ chính cậu đã là lớp keo dính hoàn hảo nhất. Và tôi hứa, tôi sẽ không để lớp keo ấy mất đi tác dụng đâu. Lớp 1-C sẽ luôn là một khối thống nhất, ngay cả khi cậu không còn ở đây.''

``Nếu cậu nói vậy thì chác tôi cũng không còn gì để lo nữa nhỉ, Ayame-san,'' tôi gật đầu, cảm nhận được sức nặng trong lời hứa của cô gái này. Rồi tôi quay sang Wata. ``Còn cậu, WataWata-san? Đừng có hơi tí là cuống cuồng lên nữa. Cậu là Phó hội trưởng đấy. Phải làm chỗ dựa vững chắc cho Inari-san và Ayame-san đấy, nghe chưa?''

``Mình hứa,'' Wata gật nhẹ, khóe mắt hơi rưng rưng. ``Mình sẽ không làm cậu thất vọng đâu, Kuon-san!''

Tôi nở một nụ cười nhẹ nhõm từ hai người, rồi từ từ chuyển ánh nhìn sang Inari. Cô nàng Lớp trưởng từ nãy đến giờ vẫn im lặng. Mái tóc trắng muốt với phần đuôi đen của giống cáo trắng khẽ rung rinh khi cô cúi gằm mặt, đôi vai nhỏ bé run lên nhè nhẹ.

``Inari-san,'' tôi gọi khẽ.

Inari ngẩng phắt đầu lên. Đôi mắt xanh lam của cô đã ngấn lệ từ lúc nào. Dù luôn mang trọng trách của một người đứng đầu lớp, nhưng tựu trung lại, Shirayuki Inari vẫn là một người rụt rè, nhút nhát, người đã từng phải nấp sau cô bạn thân vì áp lực ở những ngày đầu tiên chúng tôi đến Aogami. Nhưng giờ đây, ẩn sâu trong đôi mắt ngấn nước ấy là một ngọn lửa của sự kiên định mà tôi hiếm khi thấy.

``Kuon-san\dots'' Inari nghẹn ngào cất lời, hai bàn tay nắm chặt lấy vạt áo đến mức các khớp ngón tay trắng bệch. ``Trước khi cậu đến\dots\ mình luôn cảm thấy mình là một Lớp trưởng thất bại. Mình không nói được ai, không gắn kết được ai. Mọi người tự làm theo ý mình, còn mình chỉ biết đứng nhìn trong sự bất lực. Nhưng cậu đã cho mình thấy\dots\ một người lãnh đạo thực sự không nhất thiết phải lớn tiếng, mà là người biết cách lắng nghe và bảo vệ những người xung quanh.''

Cô bé hít một hơi thật sâu, cố gắng ngăn những giọt nước mắt trào ra, sống mũi ửng đỏ.

``Mình\dots\ mình vẫn còn rất yếu đuối. Mình vẫn hay sợ hãi. Nhưng\dots\ mình hứa!'' Inari nhấn mạnh từng chữ, giọng điệu trong trẻo vang lên đầy quyết tâm. ``Mình sẽ cố gắng trở thành một Lớp trưởng mà Kuon-san có thể tự hào! Mình sẽ bảo vệ lớp 1-C, bảo vệ thị trấn này, giống như cách cậu đã bảo vệ bọn mình vậy!''

Trái tim tôi như bị ai đó bóp nghẹn lại trước sự chân thành của cô bé. Tôi đưa tay lên, nhẹ nhàng xoa đầu Inari. Mái tóc cô mềm mại dưới lòng bàn tay tôi.

``Cậu đã là một Lớp trưởng tuyệt vời rồi, Inari-san,'' tôi mỉm cười dịu dàng. ``Đừng cố gắng trở thành bản sao của tôi hay bất kỳ ai khác. Hãy cứ là một Inari nhân hậu và chân thành. Chỉ cần như vậy, lớp 1-C tự nhiên sẽ hội tụ quanh cậu thôi.''

Inari không kìm được nữa, những giọt nước mắt nóng hổi lăn dài trên má. Cô gật đầu lia lịa, khẽ mỉm cười qua làn nước mắt. Xung quanh chúng tôi, Ayame và Wata cũng đã lặng lẽ lau khóe mi. Góc nhỏ của ban cán sự lớp giờ đây ngập tràn một sự ấm áp của sự tin tưởng và tiếp nối.

\ct

Trong khi tôi đang bận rộn dặn dò ban cán sự lớp, thì ở phía bên kia căn phòng, gần khu vực ban công lộng gió, Kaigou cũng đang bị bủa vây bởi ``phe thể thao và chiến đấu'' của lớp 1-C. Nhóm của Shiina, Yae, Akane và Kaoru đang ngồi quây tròn quanh cô nàng đuôi ngựa, ánh mắt của họ rực sáng sự kính trọng pha lẫn chút hiếu chiến đặc trưng của những người đam mê vận động.

Kaigou tựa lưng vào lan can ban công, dùng một chiếc khăn bông lau đi lớp mồ hôi trên trán sau trận ném gối. Cô nhấp một ngụm nước, đưa mắt nhìn lướt qua bốn gương mặt đang háo hức chờ đợi mình lên tiếng.

``Trông thế mà cậu mạnh đấy, Kaigou!'' Shiina đấm hai nắm tay vào nhau, đôi mắt sư tử rực lửa. ``Lúc nãy cậu ném gối bằng một tay mà lực đi căng như dây đàn! Quỹ đạo bay hoàn hảo, không có động tác thừa! Cậu có bí quyết gì không?''

Kaigou khẽ cười khẩy, xoay xoay lon nước trong tay. Dưới ánh trăng mờ ảo hắt vào từ ban công, gương mặt sắc sảo của cô trông dịu dàng hơn thường lệ, không còn vẻ phòng bị gai góc bẩm sinh.

``Bí quyết à?'' Kaigou nhướng mày. ``Chẳng có bí quyết nào ngoài việc quan sát thế giới qua lăng kính vật lý cả. Mọi chuyển động, mọi đòn tấn công của đối thủ đều bị chi phối bởi các định luật cơ học. Gia tốc, khối lượng, điểm tựa, và mô-men xoắn. Chỉ cần tính toán quỹ đạo và tìm ra điểm mù của lực tác động, cậu có thể quật ngã một kẻ to gấp đôi mình. Đó là cách sinh tồn của tôi ở cái thế giới khắc nghiệt bên kia.''

``Lăng kính vật lý\dots'' Akane lẩm bẩm, gật gù ghi nhớ như đang nghe giảng bài trên lớp. Kaoru ngồi cạnh cũng tròn mắt thán phục.

``Nhưng mà này,'' Kaigou đột ngột chuyển giọng, ánh mắt cô hạ xuống, dịu dàng một cách lạ lùng. Cô nhìn về phía những học sinh lớp 1-C đang cười đùa trong phòng. ``Suốt mùa hè ở Aogami này, mấy cậu đã dạy cho tôi một bài học đắt giá. Rằng trên đời này, tôi không thể mãi đóng cái tâm vào bản thân được. Tinh thần đồng đội, ý chí bảo vệ người khác, và\dots\ sự tin tưởng mù quáng. Ở thế giới của tôi, tin tưởng người khác là một điểm yếu chết người. Nhưng ở đây, nó lại là thứ vũ khí sắc bén nhất.''

Shiina và Yae nhìn nhau, sự tự hào ánh lên trong mắt họ.

``Bọn tớ hiểu mà,'' Yae lên tiếng, nụ cười phảng phất vẻ trưởng thành trước tuổi. Cô nàng giơ một nắm đấm lên trước mặt Kaigou. ``Dù Kaigou-san có nhìn thế giới bằng vật lý hay bằng tình cảm, thì đối với bọn tớ, cậu vẫn là một người con gái ngầu nhất. Nếu Naomi ở lại, bọn tớ hứa sẽ dùng hết sức mạnh của mình để bảo vệ con bé. Nhưng vì con bé đi cùng, nên cậu phải hứa là không được để nó khóc đâu đấy!''

``Đương nhiên rồi,'' Kaigou nhếch mép, không chần chừ mà cụng nắm đấm của mình vào tay Yae. ``Bọn tôi là gia đình của con bé mà. Ai dám đụng vào nó, tôi sẽ cho kẻ đó biết thế nào là phản lực gia tốc cực đại.''

``Đúng chuẩn Kaigou-san rồi!'' Kaoru cười lớn, vỗ tay tán thưởng.

``Nhưng mà tớ vẫn chưa cam tâm đâu nhé!'' Yae bất chợt hất mặt, chỉ tay về phía Kaigou với vẻ đầy khiêu khích. ``Hồi ở lễ hội thể thao, tớ chưa thực sự bung hết sức đâu! Kiếp này coi như không có cơ hội phục thù rồi. Cậu có dám nhận lời thách đấu ở kiếp sau không, Kaigou-san?''

Kaigou phá lên cười rạng rỡ. Đó không phải là điệu cười mỉa mai hay khinh khỉnh, mà là một nụ cười sảng khoái, giải phóng mọi tâm tư đè nén.

``Kiếp sau à?'' Kaigou đứng thẳng dậy, vuốt ngược mái tóc đen tuyền ra phía sau. ``Khá khen cho sự tự tin của cậu. Được thôi, kiếp sau hay kiếp sau nữa, bất cứ khi nào cậu tìm thấy tôi, tôi cũng sẵn sàng tiếp chiêu. Nhưng báo trước nhé, dù có đầu thai mười lần, cậu cũng không thắng nổi tôi đâu!''

``Để rồi xem!'' Shiina và Akane đồng thanh hét lên, tham gia vào màn cá cược điên rồ xuyên kiếp sống. Bọn họ cười vang, tiếng cười giòn giã hòa quyện vào gió đêm, mang theo lời hứa của những chiến binh kiêu hãnh không bao giờ chịu khuất phục trước khoảng cách của không gian và thời gian.

\ct

Bỏ lại sự ồn ào của nhóm thể thao, ở ngay trung tâm căn phòng, một nhóm khác lại mang một bầu không khí đối lập hoàn toàn: ồn ào vì\dots\ ăn uống, và đẫm nước mắt vì chia ly. Đó là ``vương quốc'' của Suzuki.

Xung quanh chiếc bàn trà lùn tịt chất đầy các loại bánh kẹo, mực xé, khoai tây chiên và những hộp bento thừa từ đêm lễ hội, Suzuki đang ngồi thu lu, tay ôm một bịch snack rong biển. Quanh cô nàng là Suzune, Honoka, Touka, Hanabi, Yuka và Saya - nhóm nữ sinh đại diện cho ``ban ẩm thực và giải trí'' của lớp 1-C.

``Cậu ăn thử cái bánh mochi đậu đỏ này đi, Suzuki-san\~{}'' Yuka nhiệt tình đẩy một hộp bánh nhỏ về phía Suzuki. ``Tôi mang từ nhà đến đấy, do chính tay bà ngoại tôi làm! Dẻo và ngọt lắm!''

``Thêm chút mực nướng đi! Nhai cho đỡ buồn miệng!'' Hanabi bồi thêm, xé một miếng mực đưa tận miệng cô nàng ngoại lai.

Suzuki há miệng nhận lấy tất cả với một tốc độ đáng kinh ngạc. Hai má cô phồng lên như sóc ngậm hạt dẻ, liên tục nhai nuốt. Trên môi cô nở một nụ cười rạng rỡ, nhưng điều kỳ lạ là, những giọt nước mắt vẫn không ngừng rơi lách tách xuống mu bàn tay. Em ấy vừa cười, vừa ăn, vừa khóc.

``Ngon quá\dots\ Oa\dots\ ngon quá đi mất\dots'' Suzuki nấc lên, giọng nghẹn ngào vì thức ăn và cảm xúc.

Ngồi ngay cạnh cô, Suzune cũng đang trong tình trạng thê thảm không kém. Cô bạn thân của Touka khóc nức nở, nước mắt nước mũi tèm lem, tay nắm chặt lấy vạt áo yukata của người bạn khác thế giới.

``S-Suzuki-chan à! Cậu đi rồi, ai sẽ ăn chung bento với tớ vào giờ nghỉ trưa đây!'' cô nàng mếu máo, khóc òa lên như một đứa trẻ. ``Ai sẽ nấu mấy món canh hầm kỳ lạ nhưng siêu ngon cho tớ ăn nữa! Tớ không muốn cậu đi đâu! Oaaaaa!''

Touka ngồi đối diện, nhẹ nhàng rút một tờ giấy ăn từ hộp ra, cẩn thận chấm những giọt nước mắt trên mặt Suzune. Đôi mắt của cô nàng Thần đồng cũng đỏ hoe, nhưng cô vẫn cố giữ vẻ điềm tĩnh thường ngày.

``Suzune-san, cậu khóc thế này thì Suzuki-chan làm sao mà yên tâm đi được,'' Touka dịu dàng nói, vỗ nhẹ lên lưng cô bạn. ``Chúng ta đã thống nhất là sẽ tiễn cậu ấy bằng nụ cười cơ mà.''

``Nhưng mà\dots\ nhưng mà\dots'' Suzune lắc đầu quầy quậy, vòng tay ôm chầm lấy Suzuki, dụi mặt vào vai cô. ``Tớ sẽ nhớ cậu ấy chết mất!''

Thấy không khí có vẻ ngày càng chìm sâu vào sự bi lụy, Honoka vội vàng vỗ tay bôm bốp để đánh lạc hướng.

``Thôi nào thôi nào! Đừng khóc nữa!'' cô nàng bờm Cáo sa mạc cố rặn ra một nụ cười tươi rói. ``Kể chuyện vui đi! Suzuki-chan, cậu còn nhớ cái hôm bọn mình tập làm kẹo bông gòn cho lễ hội không? Cái lúc mà cậu lỡ tay cho quá nhiều đường, kết quả là cái máy nó phun tơ kẹo ra trùm kín cả đầu Kuon-san ấy! Lúc đó trông anh ấy không khác gì một con người tuyết màu hồng! Hahaha!''

Câu chuyện của Honoka ngay lập tức phát huy tác dụng. Hanabi và Yuka che miệng bật cười rúc rích khi nhớ lại bộ dạng thê thảm của tôi lúc bấy giờ. Saya, người luôn kiệm lời nhất nhóm, cũng khẽ mỉm cười, ánh mắt dịu dàng nhìn Suzuki, bàn tay cô khẽ vỗ về lên lưng Suzuki như một sự an ủi lặng thầm.

Suzuki ngừng nhai, nuốt ực miếng bánh xuống cổ họng. Cô nàng bật cười thành tiếng, nhưng những giọt nước mắt lại càng tuôn rơi nhiều hơn.

``Nhớ chứ\dots\ sao mà mình quên được,'' Suzuki sụt sùi, giơ tay quệt đi dòng nước mắt đang làm nhòe tầm nhìn. Cô đưa ánh mắt trong veo nhìn một lượt những người bạn quanh bàn. ``Ở thế giới của mình\dots\ mình luôn bị coi là một kẻ kỳ dị, một gánh nặng. Mình không có bạn bè để chia sẻ đồ ăn, không có ai để cùng cười đùa những trò ngớ ngẩn. Mình luôn phải trốn tránh và chiến đấu để sinh tồn\dots''

Suzuki nghẹn giọng, bàn tay nắm chặt lấy tay Suzune và Honoka.

``Từ khi đến Aogami\dots\ các cậu đã dang rộng vòng tay đón nhận mình, ngoài Onii-sama ra. Các cậu khen món mình nấu ngon, các cậu chia cho mình những cái bánh ngon nhất, các cậu cho mình biết thế nào là mùi vị của tuổi học trò. Tớ hạnh phúc lắm. Thực sự rất hạnh phúc!''

Những lời bộc bạch từ tận đáy lòng của em ấy như một dòng suối ấm áp chảy qua trái tim mỗi người. Suzune khóc òa to hơn, nhào tới ôm chặt lấy Suzuki. Honoka không kìm được nữa, cũng đưa tay lên che mặt, nước mắt tuôn rơi ướt đẫm cả ống tay áo yukata. Hanabi, Yuka và Touka đồng loạt cúi đầu, lặng lẽ để những giọt nước mắt rơi xuống. Ngay cả Saya cũng khẽ chớp mắt, khóe mi ngân ngấn nước.

Họ không nói thêm lời nào nữa. Mọi ngôn từ lúc này đều trở nên thừa thãi. Họ chỉ cùng nhau ăn nốt những chiếc bánh cuối cùng, uống chung những ngụm trà cuối cùng, và khắc ghi hương vị của tình bạn mùa hè này sâu vào trong ký ức. Một hương vị vừa ngọt ngào của những tiếng cười, vừa mặn chát của những giọt nước mắt chia ly, nhưng chắc chắn sẽ không bao giờ phai nhạt.

\ct

Ở một góc sáng sủa nhất của căn phòng, dưới ánh đèn huỳnh quang êm dịu, Naomi đang được vây quanh bởi những ``chuyên gia sắc đẹp'' và ``tổ buôn chuyện'' khét tiếng của lớp 1-C. Saki, Miku, Mitsuki cùng với Kana, Kanade và Hotaru đang biến cô bé cựu oán linh thành một búp bê sống để thỏa sức làm dáng.

Saki đang tỉ mẩn tết lại mái tóc dài màu bạc của Naomi thành một kiểu đuôi cá cầu kỳ, trong khi Mitsuki thì vuốt ve những lọn tóc con lòa xòa trước trán cô bé.

``Này, Miku-chan, đưa mình cái kẹp tóc hình con bướm xanh kia xem nào!'' Saki vươn tay ra, ánh mắt tập trung cao độ như một nghệ nhân đang hoàn thiện tác phẩm để đời.

``Đây đây! Đừng có vội, hỏng mất nếp tóc bây giờ,'' Miku càu nhàu, nhưng tay vẫn nhanh chóng đưa chiếc kẹp nhựa đính đá lấp lánh cho cô bạn.

Kanade và Hotaru ngồi đối diện, chống cằm ngắm nhìn Naomi với ánh mắt xuýt xoa. ``Dễ thương quá đi mất\dots\ Naomi-chan mà mặc đồng phục trường mình chắc chắn sẽ trở thành hoa khôi khối năm nhất cho xem!'' Kanade cảm thán.

``Tiếc là con bé lại không thể ở lại để đi học cùng chúng ta,'' giọng Kana chợt chùng xuống, mang theo sự tiếc nuối rõ rệt. Lời nói của cô làm không khí rôm rả bỗng chốc khựng lại. Bàn tay đang tết tóc của Saki cũng dừng lại giữa không trung.

Naomi khẽ chớp đôi mắt lục bảo, nhìn những người chị gái đang vây quanh mình qua hình ảnh phản chiếu trên tấm gương nhỏ đặt trên bàn. Cô bé đưa tay lên, nhẹ nhàng chạm vào chiếc kẹp tóc hình bướm xanh vừa được cài lên tóc.

``Đẹp quá\dots'' Naomi lẩm bẩm, một nụ cười mỉm ngây thơ nhưng ẩn chứa sự sâu sắc của một linh hồn đã sống qua nhiều thế kỷ. ``Mọi người biết không\dots\ hàng ngàn năm qua, em chỉ là một linh hồn lang thang trong bóng tối của ngôi đền Aogami. Em đã quen với sự cô độc, quen với sự thù hận, và quen với việc bị con người kinh sợ. Em chưa từng nghĩ\dots\ sẽ có một ngày mình được chải tóc, được cài kẹp tóc đẹp, và được yêu thương như một con người bình thường.''

Những lời nói nhẹ bẫng của cô bé như những mũi kim đâm vào trái tim mềm yếu của các nữ sinh. Saki cắn chặt môi, cố ngăn không cho tiếng nấc bật ra. Mitsuki vội vã quay mặt đi, lấy tay quệt nước mắt.

``Cái kẹp tóc này\dots'' Mitsuki thút thít, móc từ trong túi áo ra một chiếc hộp nhỏ bằng nhung. ``Bọn chị đã hùn tiền lại mua nó ở lễ hội\dots\ định là sẽ tặng em nhân dịp em chính thức trở thành người dân của thị trấn. Ai ngờ đâu\dots\ em lại đi sớm như vậy.''

Miku lau nước mắt, cố gắng tỏ ra mạnh mẽ: ``Em giữ lấy nhé, Naomi-chan. Mỗi lần nhìn thấy con bướm xanh này, em phải nhớ đến mấy bà chị ồn ào ở thế giới này đấy. Dù em có sống bao nhiêu ngàn tuổi đi nữa, thì trong mắt bọn chị, em vẫn chỉ là một đứa em gái nhỏ cần được chăm sóc thôi!''

``Vâng ạ,'' Naomi gật đầu thật mạnh, hai tay áp chặt lấy chiếc kẹp tóc như sợ nó sẽ bay mất. Cô bé đứng dậy, dang rộng hai tay ra. ``Em có thể ôm mọi người một cái được không?''

Chẳng cần chờ đợi thêm một giây nào nữa, cả sáu cô gái đồng loạt lao tới, ôm chầm lấy Naomi thành một vòng tròn lớn. Những tiếng khóc nức nở vỡ òa, hòa lẫn với những lời dặn dò đứt quãng. ``Nhớ giữ gìn sức khỏe nhé\dots'', ``Đừng để bị bắt nạt nhé\dots'', ``Bọn chị sẽ nhớ em lắm\dots''. Naomi chìm trong vòng tay ấm áp của họ, đôi mắt xanh lục bảo lấp lánh những giọt lệ hạnh phúc. Ở thế giới này, cô bé đã thực sự tìm thấy ánh sáng cho linh hồn mình.

\ct

Trong một góc tĩnh lặng hơn, Shion đang ngồi khoanh chân, cố gắng duy trì phong thái tĩnh tại của một ``Pháp sư bóng tối'' đang thiền định. Nhưng đôi vai run rẩy và nhịp thở đứt quãng của cô nàng đã tố cáo tất cả. Ngồi cạnh cô là Aku và Sakura, cả hai đang nhìn cô nàng chuunibyou của mình bằng ánh mắt vừa buồn cười vừa thương cảm.

``Này Shion-san,'' Sakura khẽ vỗ vai cô bạn. ``Không cần phải gồng đâu. Cậu đang khóc đến mức làm nhòe cả vòng tròn ma thuật vẽ bằng bút dạ trên tay rồi kìa.''

Shion giật mình, vội vàng lấy tay áo che đi cổ tay đang lem nhem mực đen. Cô nàng hất cằm lên, giọng nói run rẩy nhưng vẫn cố tỏ ra bí hiểm: ``K-Không phải\dots\ ta không khóc! Đây là\dots\ sự cộng hưởng ma lực khi các vì sao đang đổi ngôi! Sự chia ly của Kuon-san và Kaigou-san chỉ là\dots\ là một sự dịch chuyển các dòng thời gian để duy trì sự cân bằng của vũ trụ mà thôi!''

``Vũ trụ cân bằng hay không thì tôi không biết,'' Aku cười buồn, khẽ lắc đầu. ``Nhưng tôi biết chắc là từ ngày mai, sẽ chẳng còn ai đứng ra bảo vệ bà khi lỡ miệng đâu, Shion-chan ạ.''

Câu nói của cô nàng Hầu gái như một đòn chí mạng đánh sập hoàn toàn lớp vỏ bọc chuunibyou của Shion. Cô nàng ``Pháp sư bóng tối'' cúi gầm mặt, những giọt nước mắt cuối cùng cũng phá vỡ đê bao, rơi lã chã xuống chiếu tatami.

``Ta\dots\ ta biết chứ!'' cô nàng Pháp sư nấc lên, hai tay ôm mặt. ``Ta biết là Kuon và mọi người rất quan trọng\dots\ Chỉ là ta không muốn thừa nhận là\dots\ là ta sẽ rất nhớ họ thôi!''

Cô nàng quệt nước mắt, rồi vội vã lôi từ trong túi áo ra ba chiếc bùa hộ mệnh bằng vải đỏ được khâu tay vụng về, đứng dậy chạy đến dúi vào tay tôi, Kaigou và Suzuki.

``Đây\dots\ đây không phải là bùa bình thường đâu!'' Shion cố gắng lấy lại giọng điệu kiêu ngạo, dù mắt vẫn sưng húp. ``Đây là Thánh phù Hộ thể cấp 99 do đích thân tôi truyền ma lực vào! Nó sẽ bảo vệ các ngươi khỏi mọi lời nguyền ở các thế giới khác! Hãy giữ lấy nó như một báu vật đi!''

Tôi mỉm cười nhận lấy chiếc bùa khâu bằng những đường chỉ xiêu vẹo, cảm nhận được sự ấm áp từ tấm lòng của cô. ``Cảm ơn cậu, Pháp sư Shion. Tôi hứa sẽ để nó ở vị trí trang trọng nhất trong hành trang của mình.''

Aku và Sakura khoác tay lên vai Shion, cùng nhau nở một nụ cười rạng rỡ qua những giọt nước mắt, tự hứa với nhau sẽ bảo vệ lớp học này thay cho những người sắp rời đi.

\ct

Ở phía gần cửa kéo hướng ra ban công, Nanase đang ngồi ngay ngắn trên nệm, tao nhã nâng một tách trà xanh lên môi. Phong thái điềm tĩnh, quý tộc của cô dường như không bị ảnh hưởng bởi sự hỗn loạn của bữa tiệc, nhưng nếu nhìn kỹ, sẽ thấy đôi mắt cô đang chất chứa một nỗi buồn sâu thẳm. Ngồi tĩnh lặng bên cạnh cô là Shino, người đang nhìn cô nàng với ánh mắt cảm thông.

``Là một người sinh ra trong gia tộc quyền quý,'' Nanase khẽ đặt tách trà xuống, cất giọng êm ái nhưng mang theo tiếng thở dài. ``Tôi đã quen với những cuộc chia ly. Những người đến rồi đi vì lợi ích, vì xã giao. Nhưng đây\dots\ có lẽ là lần đầu tiên, tôi cảm thấy tiếc nuối khi phải nói lời tạm biệt với những người bạn thực sự. Không toan tính, không vụ lợi.''

Shino khẽ gật đầu, hướng ánh mắt nhìn ra bầu trời đêm qua khe cửa. ``Mình cũng vậy. Ngôi đền Aogami này vốn luôn tĩnh lặng và cô đơn. Nhưng từ khi Kuon-san và mọi người đến, nó đã trở thành một nơi tràn đầy sinh khí. Sáng sớm hôm nay, mọi người sẽ bước qua cánh cổng không gian ngay tại bậc thềm của đền\dots\ Mình không thể làm gì hơn, nhưng mình sẽ tiễn họ bằng một nụ cười tươi tắn nhất, cầu chúc họ trở về nhà thượng lộ bình an.''

Shino nắm lấy tay Nanase, cả hai cùng nhìn về phía nhóm của tôi. ``Cầu chúc cho họ\dots\ dù đi đến đâu, cũng sẽ luôn tìm thấy ánh sáng và sự bình yên.''

Nanase mỉm cười, một nụ cười thanh tao nhẹ tựa lông hồng, nhưng lại chứa đựng sự trân trọng vô giá.

\ct

Trong khi mọi người đang đắm chìm trong những cuộc trò chuyện riêng, tôi tinh ý nhận ra một bóng dáng nhỏ nhắn đang đứng khoanh tay dựa lưng vào khung cửa bếp, quan sát tất cả với vẻ mặt kiêu kỳ thường thấy. Đó là Koishin.

Tôi bước lại gần, dừng lại bên cạnh cô. Cô nàg liếc nhìn tôi một cái, rồi lập tức quay mặt đi, hứ một tiếng rõ to.

``Đừng có tưởng bở,'' Koishin lên tiếng trước khi tôi kịp nói lời nào, giọng điệu sắc lẹm mang đậm chất tsundere. ``Tôi không có khóc vì anh đâu đấy nhé. Chỉ là\dots\ lông vũ từ cái gối chết tiệt nãy bay vào mắt tôi thôi.''

Tôi bật cười, không thèm vạch trần đôi mắt sưng húp của cô nàng: ``Tôi biết rồi. Bụi ở đây công nhận là nhiều thật.''

Koishin im lặng một lúc, bàn tay bấu chặt lấy gấu áo. Rồi, không thèm quay mặt lại, cô lí nhí: ``Đến thế giới mới\dots\ thì đừng có ôm rơm rặm bụng nữa. Đừng có lo chuyện bao đồng rồi tự rước họa vào thân. Giữ gìn sức khỏe đấy\dots\ đồ ngốc.''

``Tuân lệnh, cô nương,'' tôi mỉm cười nhẹ nhõm.

Lời tạm biệt gai góc ấy, với tôi, lại là lời chúc bình an chân thành nhất mà một cô nàng bằng mặt không bằng lòng có thể nói ra.

\ct

Đồng hồ chỉ về mốc bốn giờ sáng. Bầu không khí trong căn nhà đã dần chuyển từ sự náo nhiệt, ồn ào sang một vẻ tĩnh lặng đầy lưu luyến. Những câu chuyện riêng tư đã được thổ lộ, những giọt nước mắt đã được lau khô, nhưng dường như ai cũng hiểu rằng, thời khắc chia ly đang đến rất gần.

Đột nhiên, từ góc phòng, Mari rụt rè đứng dậy. Cô nàng Họa sĩ ôm khư khư một cuộn giấy lớn trước ngực, bước những bước rón rén về phía tôi, Kaigou và Suzuki. Đôi mắt hai màu của cô chớp chớp liên tục, biểu hiện rõ sự căng thẳng.

``C-Cái này\dots'' Mari ấp úng, chìa cuộn giấy ra bằng cả hai tay. ``T-Tôi và một vài bạn khác đã chuẩn bị\dots\ Tôi định lúc ra đền mới đưa, nhưng có lẽ bây giờ, tôi mới có đủ can đảm để cho các cậu xem\dots''

Tôi nhận lấy cuộn giấy, khẽ tháo sợi ruy băng màu vàng nhạt đang buộc quanh nó. Khi tôi từ từ trải rộng tờ giấy có khổ A2 ra mặt bàn, hai chị em họ cũng tò mò rướn người tới xem. Cả căn phòng lúc này đồng loạt im lặng, mọi ánh mắt đều đổ dồn về phía bức tranh.

Đó là một bức tranh vẽ tay bằng màu nước, tái hiện lại khung cảnh toàn bộ lớp 1-C đang đứng dưới cổng torii của đền Aogami trong đêm lễ hội hôm qua. Bức tranh không phải là một tác phẩm hoàn hảo, mà trông như một bức vẽ của trẻ con tiểu học hơn. Vì được vẽ gấp rút trong chỉ vài giờ đồng hồ sau khi lễ hội kết thúc, những nét cọ vẫn còn lộ rõ sự vội vã. Màu sắc ở một số chỗ chưa kịp khô đã bị lem vào nhau. Khuôn mặt của vài người thậm chí chỉ được phác họa bằng vài đường nét đơn giản.

Nhưng, chính cái sự ``chưa hoàn thiện'' ấy lại tạo nên một sức hút kỳ lạ. Bức tranh bắt trọn được linh hồn của khoảnh khắc đó: cái miệng đang cười ngoác mang tai của Shiina, vẻ mặt cau có nhưng vẫn chịu giơ tay chữ V của Koishin, sự rụt rè của Inari khi nấp sau lưng Hanabi, và cả bóng lưng vững chãi của ba chúng tôi đứng ở vị trí trung tâm, bên cạnh là Naomi đang tươi cười rạng rỡ.

``Xin lỗi mọi người\dots'' Mari cúi gầm mặt, hai ngón tay chọc chọc vào nhau. ``Tôi biết là nét vẽ còn rất lộn xộn, màu cũng lem nhem nữa. Thời gian gấp quá, tôi không kịp hoàn thiện chi tiết. Đáng lẽ tôi nên vẽ đẹp hơn để làm quà chia tay\dots''

``Cậu nói gì ngớ ngẩn vậy?'' Kaigou khoác tay lên vai Mari, khẽ bóp nhẹ. Cô nàng chỉ tay vào hình ảnh của chính mình trong tranh. ``Cậu vẽ tôi ngầu thế này cơ mà. Còn cái sự lộn xộn này\dots\ nó mới chính xác là lớp 1-C chứ!''

Tôi mỉm cười, dùng ngón tay vuốt nhẹ dọc theo viền của bức tranh như sợ làm nó nhăn nhúm. ``Đúng vậy. Nếu vẽ chỉn chu và ngăn nắp quá thì đâu còn là cái lớp cá biệt lúc nào cũng phá phách mà chúng ta biết. Sự vội vã trong từng nét cọ này, những mảng màu nhem nhuốc này\dots\ nó cho thấy cậu đã dồn biết bao nhiêu tình cảm và tâm huyết vào đây trong một thời gian ngắn ngủi.'' Tôi ngẩng lên nhìn Mari, ánh mắt đầy trân trọng. ``Cảm ơn cậu, Mari-san. Với tôi, đây đã là một kiệt tác rồi. Tôi hứa khi trở về thế giới của mình, việc đầu tiên tôi làm sẽ là mua một cái khung thật đẹp để treo nó lên tường nhà.''

Mari ngước nhìn tôi, đôi mắt vàng-đỏ mờ đi vì sương lệ. Cô cười tươi rạng rỡ, gật đầu liên tục thay cho lời nói.

Ngay khi Mari vừa lùi lại, Uzuki - cô nàng nhiếp ảnh gia năng nổ của lớp - vội vã chạy tới. Trên tay cô là một cuốn album nhỏ xíu, vẫn còn thoang thoảng mùi hóa chất tráng phim ngai ngái.

``Mình có thứ này muốn tặn các cậu!'' Uzuki nhét cuốn album vào tay tôi, thở hổn hển. ``Mình đã tranh thủ mượn phòng tối tạm thời ở tiệm nhiếp ảnh thị trấn để rửa nhanh những bức ảnh chúng ta chụp ở lễ hội tối qua. Tuy màu sắc chưa được sắc nét lắm, nhưng\dots\ nhưng đây là những khoảnh khắc chân thực nhất.''

Tôi lật mở những trang album mỏng manh. Trong đó là những bức ảnh đen trắng và ngả vàng hoài cổ: cảnh lớp 1-C hì hục dựng gian hàng, cảnh Kaigou lườm đám nam sinh trường khác, cảnh Suzuki đang ôm bụng cười ngặt nghẽo bên đĩa takoyaki, và cảnh Naomi đang ngước nhìn pháo hoa với đôi mắt sáng rực. Những khoảnh khắc tưởng chừng giản đơn nhưng giờ đây lại nặng trĩu ý nghĩa.

Đúng lúc bầu không khí đang chìm trong sự hoài niệm lắng đọng, chiếc điện thoại trong túi áo yukata của tôi bỗng rung lên bần bật. Màn hình tự động sáng rực, phóng ra một luồng sáng ảo ảnh mờ nhạt giữa không trung. Và rồi, một giọng nói máy móc vang lên đủ lớn để cả phòng cùng nghe thấy.

\eng{Tiến hành xuất báo cáo thống kê cuối cùng của Hệ thống tại không gian Aogami.}

Game bỗng dưng tự kích hoạt hệ thống loa ngoài, khiến tất cả mọi người đều giật mình ngơ ngác.

\eng{[Tổng kết dữ liệu hành vi:]}\\
\eng{- Số lần Kuon-user thở dài trong suốt thời gian ở Aogami: 847 lần.}\\

\eng{- Số lần Kaigou-user đe dọa bẻ cổ người khác: 312 lần.}\\
\eng{- Số lần Suzuki-user mỉm cười: Vượt quá giới hạn bộ đếm.}\\
\eng{- Số lần Naomi-user hỏi mọi người bằng sự tò mò vô tận: Vượt quá giới hạn bộ đếm.}\\

Màn hình ảo ảnh nhấp nháy vài giây trước khi Game chốt lại bằng một giọng điệu dường như có chút mỉa mai hơn.

\eng{- Số lần Nanase và Sakura cãi nhau chí chóe vì những chuyện cỏn con: 32 lần.}\\
\eng{- Số lần Lớp trưởng Shion tự làm trò cười vì những câu thần chú ngớ ngẩn bị lỗi nhịp: 402 lần.}

``C-Cái gì mà làm trò cười hả cái hệ thống ngốc nghếch kia! Đó là do nhiễu loạn ma lực!'' Shion đỏ mặt tía tai gào lên, vung vẩy hai tay phản đối.

Sakura và Nanase thì đồng loạt liếc xéo nhau, cùng hứ một tiếng rồi quay ngoắt mặt đi chỗ khác. Còn Kaigou thì nhướng mày: ``Chỉ có 32 lần thôi sao? Chắc là tôi tính thiếu lúc ngủ rồi.'

\eng{[Kết luận: Đây là vũ trụ mà Hệ thống đã quan sát được nhiều biến số cảm xúc nhất. Đánh giá tổng thể: Tuyệt vời.]}

Một sự im lặng ngắn ngủi bao trùm lấy căn phòng, trước khi tất cả vỡ òa. Cả lớp bật cười phá lên. Những tiếng cười giòn giã hòa lẫn với những giọt nước mắt vẫn còn đọng trên khóe mi. Một ``bản báo cáo'' khô khan của máy móc nhưng lại là sự ghi nhận chân thực nhất về những tháng ngày chúng tôi đã cùng nhau trải qua. Thậm chí một AI vô tri như Game dường như cũng bị tình người nơi đây làm cho rung động.

Đợi cho tràng cười lắng xuống, từ phía ban công, Yuki và Azuki cùng nhau bước tới. Cô nàng tóc đen-đỏ đang cầm trên tay một chiếc USB nhỏ màu bạc, còn cô nàng có mái tóc đen-tím với mái tóc cắt bob đi bên cạnh, trên môi nở một nụ cười hiền dịu.

``Bọn mìnnh cũng có quà,'' Yuki lên tiếng, giọng nói trầm ấm thường ngày nay có pha chút nghèn nghẹn. Cô nàng chìa chiếc USB ra trước mặt tôi. ``Không có vật chất gì to tát đâu. Chỉ là một bài hát do Azuki sáng tác lời, còn mình soạn nhạc thôi. Hai chúng mình đã thu âm nó ngay tại phòng phát thanh của trường chiều hôm qua.''

``Bọn mình định hát trực tiếp cho mọi người nghe\dots'' Azuki khẽ nói, hai tay đan vào nhau. ``Nhưng mà\dots\ bọn mình sợ đang hát giữa chừng lại khóc mất, nên đành thu âm lại. Kuon-san, cậu có thể bật nó lên bây giờ được không?''

``Tất nhiên rồi,'' tôi gật đầu, vội vàng lấy chiếc laptop của mình ra, cắm USB vào và kết nối với chiếc loa nhỏ đặt ở góc phòng.

Một tệp âm thanh duy nhất nằm trơ trọi trong đó, mang tên: ``Mùa hè cuối cùng ở Aogami''. Tôi nhấp chuột ấn nút ``Play''.

Vài tiếng lạch cạch của thiết bị thu âm vang lên, tiếp theo là một tràng ho khan hắng giọng của Yuki. Và rồi, những nốt nhạc guitar acoustic mộc mạc, trong trẻo bắt đầu vang lên. Không có sự hỗ trợ của bất kỳ nhạc cụ điện tử nào, chỉ có tiếng dây đàn gảy từng nhịp chậm rãi, mang theo âm hưởng của gió chiều và tiếng ve sầu mùa hạ.

Giọng hát thanh thoát, bay bổng của Azuki cất lên, hòa quyện hoàn hảo với tiếng đệm guitar của Yuki:

\begin{center}
  ``Một chiều tháng Bảy nắng đổ trên thềm hiên,\\
  Cơn gió lạ thổi đến từ một thế giới bình yên.\\
  Mang theo những nụ cười, xua đi những muộn phiền,\\
  Đánh thức những trái tim vốn đã ngủ quên\dots''
\end{center}

Cả căn phòng chìm vào tĩnh lặng tuyệt đối. Mọi người đều nhắm mắt lại, đắm chìm trong giai điệu nhẹ nhàng ấy. Từng câu chữ do Azuki viết không hề hoa mỹ, chỉ là những hình ảnh chân thực nhất về mùa hè mà chúng tôi đã trải qua. Về những buổi chiều nằm vắt vẻo trên sân thượng, về tiếng cười giòn giã trong lễ hội thể thao, về những giọt mồ hôi rơi trên đường đua, và về đêm pháo hoa rực rỡ ở đền Aogami.

\begin{center}
  ``Dù ngày mai đây phương trời mỗi người mỗi ngả,\\
  Dòng thời gian trôi, ký ức xin đừng phôi pha.\\
  Cảm ơn cơn gió mùa hạ đã ghé qua,\\
  Cho ta biết thế nào là một mái nhà\dots''
\end{center}

Khi những nốt nhạc cuối cùng từ từ nhỏ dần và kết thúc bằng một tiếng ngân dài của guitar, không gian trong phòng vẫn giữ nguyên sự vắng lặng. Chỉ còn tiếng ve kêu ngoài hiên và tiếng thở khe khẽ của mọi người.

Tôi ngồi lặng người đi. Đôi mắt tôi dán chặt vào chiếc màn hình laptop đã chuyển sang màu đen. Một cảm giác bồi hồi chực trào nơi cuống họng. Ở thế giới gốc, tôi đã từng nghe vô số bài hát hay, được sản xuất bởi những dàn âm thanh hiện đại bậc nhất, được biểu diễn ở những sân khấu hoành tráng. Thế nhưng, bản thu âm thô mộc, còn lẫn cả tiếng lật giấy sột soạt này lại chạm đến những ngóc ngách sâu thẳm nhất trong tâm hồn tôi.

Bức tranh lộn xộn của Mari, và bản nhạc acoustic mộc mạc của Yuki và Azuki. Chúng không phải là những món quà đắt tiền. Chúng là những mảnh ghép ký ức cuối cùng, là kết tinh tình cảm vô giá mà lớp 1-C dành trọn cho chúng tôi. Một sự trân trọng tuyệt đối dâng lên trong lồng ngực tôi. Tôi đã đến thế giới này với thân phận một kẻ ngoại lai, mang theo những vết thương rỉ máu từ một vũ trụ khác. Nhưng chính nơi đây, chính những con người này đã dùng tình yêu thương nguyên bản nhất để chữa lành tôi, và để tôi có cơ hội được chữa lành họ.

Tôi từ từ đứng dậy, hướng về phía Yuki, Azuki, Mari, và toàn thể lớp 1-C, những người bây giờ đang chìm trong giấc ngủ sau nhiều tiếng thức đêm.

Tôi quay sang nhìn ba người đồng hành của mình. Kaigou đang khoanh tay dựa lưng vào tường, Suzuki thì ôm chặt chiếc balo đựng đầy những hộp bento rỗng mà cô nhóc quyết định mang theo làm kỷ niệm, còn Naomi đang nắm chặt vạt áo của Kaigou.

``\vie{Nhìn mọi người ngủ say sưa thế kia,}'' Kaigou cất giọng thì thầm, khóe môi khẽ nhếch lên tạo thành một nụ cười ấm áp hiếm hoi. ``\vie{Đúng là tuổi trẻ có khác. Quậy phá cho đã rồi lại gục ngã không biết trời đất gì.}''

``\vie{Em sẽ nhớ mọi người ở đây lắm,}'' Suzuki sụt sùi, khẽ lau khóe mắt. ``\vie{Em chưa bao giờ có nhiều bạn bè như thế này cả\dots}''

Naomi ngước đôi mắt lục bảo lên nhìn tôi, giọng nói trong veo mang theo một chút ngập ngừng: ``\vie{Onii-chan\dots\ chúng ta thực sự phải đi sao?}''

Tôi tiến lại gần, đặt tay lên mái tóc bạc của cô bé, gật đầu chậm rãi. ``\vie{Ừ. Thế giới này vốn dĩ đã có quỹ đạo riêng của nó. Sự xuất hiện của chúng ta\dots\ có lẽ cũng sẽ không cần thiết nữa. Đã đến lúc phải trả lại sự yên bình vốn có cho nơi này rồi.}'' Tôi nhìn lướt qua Kaigou và Suzuki. ``\vie{Mọi người đã sẵn sàng chưa?}''

``\vie{Lúc nào cũng sẵn sàng,}'' Kaigou đấm nhẹ hai tay vào nhau.

Tôi hít một hơi thật sâu, rồi lại quay mặt về phía những học sinh đang say giấc. Tôi cúi gập người xuống, một cái cúi chào thấp nhất, chân thành nhất.

``Cảm ơn mọi người,'' giọng tôi khàn đi, thì thầm trong tĩnh lặng để không đánh thức họ. ``Đây là những món quà tuyệt vời nhất mà tôi từng nhận được. Hikawa Kuon này, dù có đi đến tận cùng vũ trụ, cũng sẽ không bao giờ quên mùa hè ở Aogami, không bao giờ quên lớp 1-C chúng ta.''

Kaigou, Suzuki và Naomi cũng lần lượt đứng lên, xếp hàng ngang bên cạnh tôi và cúi gập người.

``Cảm ơn mọi người vì tất cả!'' bốn chúng tôi đồng thanh thì thầm, để mặc cho những giọt nước mắt cuối cùng của sự tri ân tự do rơi xuống.

Bên ngoài cửa sổ, màn đêm mờ ảo đã bắt đầu nhường chỗ cho một sắc xanh xám của rạng đông. Những tia sáng đầu tiên của ngày mới đang từ từ nhô lên phía sau rặng núi xa xa, hắt những vệt sáng nhạt nhòa vào căn phòng.

Đã đến lúc rồi. Thời khắc của sự bắt đầu mới, và cũng là thời khắc để nói lời chào tạm biệt.

Bữa tiệc ngủ độc nhất vô nhị của chúng tôi đã dần đi tới những giờ khắc cuối cùng.

Đúng lúc này, chiếc đồng hồ trên cổ tay tôi rung nhẹ. Giọng nói quen thuộc của Game vang lên trong đầu, xua đi bầu không khí tĩnh mịch.

\eng{Giờ chuẩn bị đã điểm rồi đó, Kuon-user. Lúc này đã là 4:50 sáng. Vòng xoáy không gian đang chuẩn bị đạt trạng thái ổn định nhất.}

Tôi hít một hơi thật sâu rồi vỗ tay mấy cái thật to. ``\textbf{Dậy thôi mọi người ơi! Sáng rồi!}''

Kaigou, Suzuki và Naomi cũng hùa theo, đi lay từng người dậy. Lớp 1-C lục tục tỉnh giấc, những tiếng ngáp ngắn ngáp dài vang lên khắp phòng. Thế nhưng, khác với sự uể oải, ngay khi nhớ ra hôm nay là ngày gì, tất cả đều bật dậy và bắt đầu dọn dẹp mọi thứ một cách nhanh gọn đến đáng kinh ngạc. Những chiếc gối vương vãi lông chim được gom lại, chăn nệm được xếp vuông vức, và rác từ bữa tiệc đêm qua nhanh chóng được phân loại cho vào túi.

Trong lúc lớp 1-C đang dọn dẹp phòng khách, gia đình Hikawa chúng tôi lui vào phòng trong để thu dọn những hành trang cuối cùng. Chúng tôi cẩn thận nhét vào balo những món đồ lưu niệm nho nhỏ, và đặc biệt là những bộ quần áo bình thường mà chúng tôi đã mua cho Naomi trong ngày đầu tiên em ấy tái sinh tại thời đại này. Xong xuôi, ba chúng tôi thay lại bộ trang phục chiến đấu quen thuộc mang sắc đen tối màu mà chúng tôi vẫn thường mặc ở vũ trụ cũ, cẩn thận gấp gọn gàng những bộ yukata sặc sỡ mượn từ lễ hội tối qua và đặt ngay ngắn trên nệm.

Khi mọi người đã tề tựu đông đủ ở cửa ra vào, tôi nán lại vài giây, xoay người nhìn bao quát căn hộ nhỏ bé này một lần cuối cùng. Ánh sáng nhạt nhòa của buổi rạng đông hắt qua ô cửa sổ, chiếu lên mặt bàn gỗ, chiếc tivi cũ, và gian bếp nhỏ nơi chúng tôi đã có những bữa ăn rộn rã tiếng cười. Căn nhà này không quá rộng lớn, cũng chẳng sang trọng, nhưng nó đã chứng kiến trọn vẹn những kỷ niệm vui buồn, đã bao bọc và sưởi ấm tâm hồn chúng tôi suốt năm tháng ròng rã vừa qua. Nơi đây, thực sự đã từng là một mái ấm.

Tôi khẽ mỉm cười, thì thầm một tiếng ``Cảm ơn vì đã cưu mang chúng tôi,'' rồi dứt khoát bước qua bậc cửa.

\ct

Năm giờ sáng.

Trời tờ mờ sáng. Sương mù vẫn còn đọng lại trên những mái nhà lợp ngói âm dương cổ kính của thị trấn Aogami. Bầu trời bình minh nhuộm màu tím hồng, không khí trong lành sau cơn mưa sao băng đêm qua. Hơi lạnh của buổi sớm mai phả vào mặt, mang theo mùi hương ngai ngái của đất ẩm và lá cây. Không gian tĩnh mịch đến mức chúng tôi có thể nghe rõ tiếng bước chân của hàng chục người vang vọng trên con đường đá lát.

Đoàn người kéo dài từ căn hộ nhà Hikawa tiến dần về phía ngọn núi nơi đền Aogami tọa lạc. Kaigou dắt tay Naomi đi ở phía trước, Suzuki đi cạnh Suzune. Toàn bộ lớp 1-C đi theo sau trong một sự im lặng kỳ lạ, với cô chủ nhiệm Yuko-sensei đi sau cùng. Đó không phải là sự im lặng của nỗi buồn bi lụy, mà là một sự ``im lặng chấp nhận''. Họ đã khóc đủ rồi. Giờ đây, họ đang dùng những bước chân kiên cường nhất để tiễn những người bạn của mình đi đến cuối con đường.

Đi dạo cùng bước chân đều đặn của mọi người, tâm trí tôi bất giác trôi dạt về những suy nghĩ miên man.

Dọc theo lộ trình hướng về phía núi, thị trấn Aogami hiện lên qua từng lăng kính của kỷ niệm. Chúng tôi đi ngang qua khu phố chính, nơi những dải đèn lồng đầy màu sắc và các khung gỗ của gian hàng lễ hội tối qua vẫn còn nằm lặng im chờ được thu dọn.

Bước qua ngã tư, nhà ga Aogami hiện ra với kiến trúc mái vòm quen thuộc, nơi ``vòng lặp mất tích bí ẩn'' từng giam cầm thị trấn trong nỗi sợ hãi giờ đây đã hoàn toàn chìm vào dĩ vãng.

Đoàn người tiếp tục rảo bước qua khu trung tâm thương mại với đài phun nước bằng đá cẩm thạch, nơi mà tuần đầu tiên chúng tôi được trải nghiệm một ngày nghỉ hiếm hoi.

Chúng tôi lẳng lặng đi dọc theo con đường mòn bên bờ ven sông róc rách nước chảy. Cuối cùng, hình bóng trường Trung học Aogami hiện ra ở phía xa xa, yên tĩnh dưới bầu trời đang chuyển dần sang sắc vàng. Những địa điểm ấy đã gói ghém trọn vẹn mọi kỷ niệm buồn vui của chúng tôi trong suốt năm tháng qua.

Thu lại ánh nhìn từ khung cảnh thị trấn, tôi đảo mắt quan sát một lượt những gương mặt thân quen của lớp 1-C đang rảo bước bên cạnh. Một tập thể hoàn hảo, một bức tranh trọn vẹn của tuổi thanh xuân.

Thế nhưng, sự trọn vẹn này lại khiến tôi cảm thấy rùng mình.

``Này, Game,'' tôi gọi thầm.

``\eng{Gì vậy, Kuon-user?}'' giọng nói máy móc quen thuộc vang lên khe khẽ từ chiếc đồng hồ trên tay.

``\eng{Ông có để ý không?}'' tôi rảo bước, đôi mắt lướt qua Azuki đang đi cạnh Yuki. ``\eng{Azuki\dots\ ở thế giới cũ của tôi, cô ấy chưa lên văn phòng  vào thời điểm đó. Nhưng ở đây, cô ấy lại hiện diện như một phần không thể thiếu. Và ngược lại\dots}''

Khóe mắt tôi giật nhẹ khi nhận ra một khoảng trống vô hình trong đội hình của lớp. Một sự vắng mặt mà suốt cả mùa hè này tôi đã không kịp nhận ra vì quá bận rộn.

``\eng{Rushia-chan. Pháp sư gọi hồn. Cô ấy không hề có một dị bản nào tồn tại ở lớp 1-C này cả.''}

Ông ta ``hừm'' một tiếng dài, như đang xử lý một lượng dữ liệu khổng lồ trước khi trả lời.

``\eng{Dòng chảy của các vũ trụ song song luôn có những sự giao thoa với những đặc điểm tinh tế. Thế giới Aogami này là một hệ sinh thái được tạo ra từ những mảnh vỡ linh hồn và tiềm thức mạnh mẽ nhất. Việc một cá thể xuất hiện sớm hơn dự kiến, như trường hợp của Azuki, là một điềm báo tốt lành cho sự hòa nhập của cô ấy ở vũ trụ gốc trong tương lai gần.}''

``\eng{Tuy nhiên\dots\ việc một cá thể hoàn toàn vắng mặt ở một thế giới thanh bình như thế này\dots\ cậu có thể hiểu đó là một sự đào thải của định mệnh. Linh hồn của Uruha Rushia ở vũ trụ gốc đang dần mất đi sợi dây liên kết với thế giới đó. Sự vắng mặt của cô ấy ở đây chính là một lời tiên tri tàn nhẫn cho sự đứt gãy sắp tới ở thế giới của cậu.}''

Sống lưng tôi lạnh toát. Tôi đã biết trước sự tàn nhẫn của thế giới cũ, nhưng khi thấy nó được phản chiếu qua chính sự thanh bình của Aogami, tôi mới cảm nhận được trọn vẹn sức nặng của lịch sử. Dòng chảy định mệnh vẫn đang âm thầm vận hành, tàn nhẫn và không thể đảo ngược.

``\textit{Ra là vậy à\dots}''

Tôi khép hờ đôi mắt, hít một hơi sương sớm vào lồng ngực để xua tan đi sự nặng nề.

\ct

Chẳng mấy chốc, cả đoàn đã đến dưới chân những bậc thang đá rêu phong dẫn lên đền Aogami. Những cánh cổng torii màu đỏ son đứng sừng sững trong sương sớm như những người lính gác thầm lặng.

Đoàn người chầm chậm bước lên những bậc thang đá dốc đứng phủ đầy rêu phong. Tiếng bước chân đều đặn vang lên giữa không gian tĩnh mịch. Khi đi được nửa đoạn đường, Shino bỗng dừng lại ở một bậc thềm nhô rara, chính là nơi cô đã từng chắp tay cầu nguyện thần linh cứu rỗi thị trấn trong vòng lặp quá khứ.

Cô nàng tóc nâu quay đầu nhìn xuống cả lớp, rồi hướng ánh mắt về phía ba chúng tôi. Một nụ cười nhẹ nhõm và thanh thản nở trên môi cô. ``Lời cầu nguyện năm ấy\dots\ cuối cùng cũng được đáp lại. Dù chỉ trong một mùa hè ngắn ngủi,'' Shino thì thầm, đôi mắt ánh lên niềm hạnh phúc trọn vẹn trước khi tiếp tục bước lên.

Tại đỉnh của bậc thang, Sakura đã đứng chờ sẵn tự bao giờ. Cô mặc bộ trang phục miko truyền thống, mái tóc bạc buộc gọn gàng bằng sợi chỉ đỏ. Dưới ánh sáng mờ ảo của buổi bình minh, cô nàng trông như một nữ thần canh giữ ranh giới giữa nhân gian và cõi linh thiêng.

``Mọi người đến rồi,'' Sakura cất giọng êm ái, đi tới đón chúng tôi. Cô chắp tay lại trước ngực, đôi mắt nhắm nghiền, bắt đầu lẩm nhẩm những lời cầu nguyện cổ xưa bằng một thứ ngôn ngữ du dương.

Mỗi lời cầu nguyện của cô như hóa thành những đốm sáng nhỏ li ti bay lơ lửng trong không trung, hòa vào màn sương sớm, mang theo lời chúc bình an và sự bảo hộ tuyệt đối từ thần linh của đền Aogami.

``Ngôi đền này đã chứng kiến sự khởi đầu của mọi thứ,'' Sakura mở mắt ra, nhìn thẳng vào tôi. ``Và giờ, nó sẽ chứng kiến sự kết thúc. Cho dù các cậu có đi xa tới đâu, những lời cầu nguyện tới những điều tốt đẹp nhất của chúng tôi, và của cả thị trấn Aogami này\dots\ sẽ mãi luôn theo các cậu.''

Tiếp bước Sakura, toàn bộ lớp 1-C cùng với Yuko-sensei đều đồng loạt chắp tay lại. Tiếng vỗ tay vang lên hai nhịp dứt khoát, phá vỡ sự tĩnh lặng của sớm mai. Tất cả đều nhắm nghiền mắt, cúi đầu thành kính. Không một ai nói ra, nhưng tôi biết, họ không cầu xin thần linh cho bản thân, mà họ đang dồn tất cả tâm niệm để cầu chúc cho chuyến hành trình vượt không gian của chúng tôi được bình an vô sự. Một luồng gió nhẹ lướt qua, làm tiếng chuông đền ngân vang một hồi dài, mang theo những lời nguyện ước ấy bay thẳng lên bầu trời đang rạng sáng.

Tiếng chuông vừa dứt, ngay phía sau lưng chúng tôi, giữa khoảng sân rộng lớn của đền, không khí đang vặn xoắn lại một cách dữ dội. Một lỗ xoáy khổng lồ mang màu sắc trắng đen đan xen đang từ từ mở ra. Nó không mang màu tím u ám và thù địch như những cổng không gian trước đây, mà lại phát ra một vầng hào quang nhẹ nhàng, thanh bình, giống hệt như âm dương hòa hợp.

\img{portal}

Đó chính là con đường đưa chúng tôi rời khỏi thế giới này. Con đường duy nhất để có cơ hội trở về nhà.

Tôi quay lại đối mặt với lớp 1-C. Tất cả mọi người đã đứng thành một hàng ngang, không ai bảo ai, những giọt nước mắt lại một lần nữa chực trào.

``Cảm ơn cậu, Shino-san,'' tôi bước tới, khẽ nắm lấy tay cô nàng miko tóc trắng. ``Vì đã cho chúng tôi một kỷ niệm tuổi trẻ không thể nào quên.''

Shino lắc đầu, nước mắt lăn dài trên má: ``Chính các cậu mới là những người đã cứu chúng mình đó. Xin hãy bảo trọng, Kuon-san.''

Bên cạnh tôi, Kaigou tiến đến trước mặt Shiina và Yae. Không có những lời sướt mướt hay những cái ôm yếu đuối. Kaigou chỉ nhếch mép cười, gật đầu một cái thật mạnh, và giơ nắm đấm ra. Shiina và Yae cũng đồng loạt mỉm cười, vung tay cụng những nắm đấm của mình vào tay Kaigou. Không cần lời nói nào cũng hiểu được sự trân trọng họ dành cho nhau.

Trong khi đó, Suzuki đang ôm chặt lấy Suzune và Touka. Cô em gái của Kaogou nức nở đến mức không thở nổi, vùi mặt vào hõm vai của hai người bạn thân thiết nhất. ``Hức\dots\ mình sẽ không bao giờ quên mọi người đâu\dots'' Suzuki ngẩng lên, cố gắng nặn ra một nụ cười rạng rỡ nhất qua đôi mắt đẫm lệ. ``Hẹn gặp lại ở một thế giới nào đó nhé!''

Naomi lùi lại một bước, hướng về phía toàn thể lớp 1-C. Cô bé nắm chặt lấy chiếc kẹp tóc hình bướm xanh trên đầu, rồi từ từ gập người xuống sâu nhất có thể. ``Cảm ơn mọi người vì đã cho em một cuộc sống mới!''

Khi Naomi vừa ngẩng lên, Yuko-sensei bước lên phía trước. Đôi mắt cô giáo đỏ hoe, cô chỉ nói một câu giản đơn: ``Sống khỏe nhé, Naomi-chan.'' Rồi, cô lặng lẽ đứng phía sau học sinh gật đầu, rồi ra hiệu cho lớp trưởng.

Từ phía sau, Nanase bước tới, trao cho tôi một phong thư dày cộp, được dán kín bằng sáp đỏ.

``Đây là bức thư viết tay của cả lớp,'' cô nàng mỉm cười dịu dàng, đôi mắt ánh lên sự tự hào. ``Khi nào mọi người và mọi người trở về nhà an toàn, hãy mở nó ra nhé.''

Tôi cẩn thận cất bức thư vào túi áo trước ngực, áp sát vào trái tim mình: ``Tôi hứa. Nhất định sẽ đọc nó.''

Nanase gật đầu, rồi chuyển sang món đồ thứ hai. Đó là một bộ đồng phục trường Trung học Aogami được gấp gọn gàng, bên trên là những bộ yukata sặc sỡ mà chúng tôi đã mặc tối qua. Cô trao chúng cho bốn người chúng tôi.

``Còn đây là đồng phục chính thức của trường Trung học Aogami,'' Nanase nói, giọng có chút run rẩy. ``Dù mọi người có đi đến chiều không gian nào, các cậu vẫn sẽ luôn là học sinh của lớp 1-C. Còn những bộ yukata này\dots\ xin mọi người đừng trả lại. Hãy mang nó theo làm kỷ vật của mùa hè này.''

Hai người họ cắn chặt môi để không bật khóc, gật đầu nhận lấy những kỷ vật vô giá. Phía sau lưng, lỗ xoáy không gian trắng đen bắt đầu xoay nhanh hơn, tỏa ra một lực hút nhẹ nhàng nhưng hối thúc.

Tôi nhìn chằm chằm vào lỗ xoáy khổng lồ mang hai nửa đen trắng đang cuộn trào trước mặt, cảm nhận được một luồng sức mạnh vô hình đang nhẹ nhàng kéo lấy cơ thể.

``Game,'' tôi gọi thầm từ chiếc đồng hồ, ``\eng{Đây có thực sự là cánh cổng đưa chúng ta về nhà không?}''

``\eng{Không thể chắc chắn 100\% được, Kuon-user,}'' giọng nói của Hệ thống đáp lại, có chút rè rè do nhiễu loạn không gian. ``\eng{Nhưng chỉ còn một cách duy nhất để biết thôi.}''

Từ bên cạnh tôi, Naomi hơi nép người lại, ánh mắt rụt rè nhìn vào trung tâm của lỗ xoáy. ``Ông Game này\dots'' cô bé lẩm nhẩm. ``Nếu em bước qua đó\dots\ liệu em có thể giữ được sự toàn vẹn của mình không? Em có bị tan biến không?''

Game khẽ ngập ngừng một nhịp. ``\eng{Khó nói lắm, Naomi. Cơ thể hiện tại của cô được cấu thành từ linh lực của đền Aogami. Việc bước ra khỏi không gian này tiềm ẩn một rủi ro cực kỳ lớn. Nhưng\dots\ nếu cô thực sự quyết tâm muốn đi cùng họ đến như thế\dots\ sao không thử một phen xem sao?}''

Naomi chớp mắt, rồi nở một nụ cười rạng rỡ, bàn tay nhỏ bé siết chặt lấy tay tôi.

Tôi quay lại phía sau, nhìn tập thể lớp 1-C đang nín thở theo dõi. Tôi, Kaigou, Suzuki và Naomi cùng nhau cúi chào một lần cuối cùng.

Và rồi, từ chiếc đồng hồ trên tay tôi, âm thanh văng vẳng của Game phát ra, đủ lớn để cả lớp cùng nghe thấy: ``Tôi rất vinh dự vì đã được đồng hành với tất cả mọi người. Mọi người sống khỏe nhé.''

Tất cả mọi người đồng thanh hô lên: ``Mọi người cũng vậy nhé.''

``Được rồi,'' tôi hít một hơi thật sâu, dang rộng hai tay ra để Kaigou, Suzuki và Naomi cùng nắm lấy nhau thành một vòng tròn khép kín. ``Cùng đếm nhé.''

``\textbf{Một!}'' Kaigou hô lớn.

``\textbf{Hai!}'' Suzuki nhắm tịt mắt lại.

``\textbf{Ba!}'' Naomi dõng dạc hét lên.

Ngay khi con số ba vừa dứt, một lực hút cực mạnh bất ngờ bùng nổ. Cơ thể chúng tôi nhẹ bẫng đi, nhấc bổng khỏi mặt đất và lập tức bị lỗ xoáy hút tọt vào trong.

Bên trong đường hầm không gian, mọi thứ quay cuồng với tốc độ chóng mặt. Những vệt sáng đen trắng đan xen xé toạc võng mạc.

``\textbf{Bám chặt lấy nhau! Tuyệt đối không được buông tay!}'' tôi gào lên giữa tiếng gió rít gầm gào.

Suzuki hét lên thất thanh, ôm chầm lấy Kaigou - người chị gái của mình - chặt đến mức các khớp ngón tay trắng bệch. Ở phía đối diện, Naomi thu mình lại, vùi đầu ôm chặt lấy lồng ngực tôi. Lỗ xoáy bắt đầu nhấp nháy một thứ ánh sáng trắng chói lòa, tốc độ cuộn xoáy mỗi lúc một nhanh hơn, nhanh đến mức ý thức của tôi bắt đầu mờ dần đi\dots

\pov{Misora Shino}

Bóng dáng của bốn người Kuon vừa bị hút vào trong lỗ xoáy thì một luồng ánh sáng trắng chói lọi bùng lên lóa mắt.

*BỤP!*

Tôi đưa tay lên che mặt. Khi ánh sáng ấy dịu đi, lỗ xoáy đen trắng khổng lồ cũng đã co rúm lại và biến mất hoàn toàn vào hư không, không để lại dù chỉ là một tiếng động tĩnh điện nhỏ nhất. Cánh cổng không gian đã chính thức khép lại.

Khuôn viên đền Aogami trở lại vẻ tĩnh mịch thường ngày. Cánh cổng torii màu đỏ son vẫn đứng sừng sững, không có gì thay đổi, chỉ có những tia nắng vàng ươm của buổi sớm mai đang xuyên qua kẽ lá, chiếu rọi xuống khoảng sân trống trải.

Không còn Kuon. Không còn Kaigou, Suzuki. Và cũng không còn bóng dáng của oán linh Naomi.

Một khoảng không trống rỗng đột ngột bao trùm lấy tôi. Không phải là sự trống trải của tuyệt vọng, mà là một cảm giác mất mát dịu dàng, chênh vênh đến kỳ lạ. Cứ như thể một mảnh ghép quan trọng nhất của mùa hè này đã bị tước đi mãi mãi. Tôi từ từ hạ tay xuống, hai chân bỗng chốc mềm nhũn ra, vô thức quỳ gối xuống nơi bốn người họ vừa đứng cách đây vài giây.

Tôi áp hai tay lên ngực trái, nhắm mắt lại, lắng nghe nhịp đập đều đặn của trái tim mình. Nhịp đập của sự sống. Nhịp đập của một tương lai bình yên mà họ đã dùng cả mạng sống để mang đến cho thị trấn này. Khóe mắt tôi cay xè, những giọt nước mắt ấm nóng lăn dài trên má, rớt xuống mặt đất đá lạnh lẽo.

Dưới ánh nắng bình minh ấm áp, những cánh hoa anh đào màu hồng phấn đang từ từ rơi lất phất trong gió\dots\ dù cho lúc này, vẫn đang là giữa mùa hè. Phép màu cuối cùng của đền Aogami đang nhẹ nhàng vỗ về nỗi buồn chia ly của lớp 1-C.

``Cảm ơn mọi người\dots\ và tạm biệt,'' tôi mỉm cười qua những giọt nước mắt, khẽ thì thầm với khoảng không vô tận.

\begin{center}
  \uline{Nhiều ngày sau.}
\end{center}

Một tuần lễ đã trôi qua kể từ buổi sáng kỳ diệu ở đền Aogami.

Thị trấn của chúng tôi đang đắm mình trong những ngày rực rỡ nhất của tháng Tám. Tiếng ve sầu râm ran vang dội từ những tán cây cổ thụ dọc theo con đường dẫn tới nhà ga trung tâm. Tôi bước đi chậm rãi trên vỉa hè, tận hưởng cái nắng vàng ươm rớt xuống qua kẽ lá, một điều tưởng chừng như quá đỗi bình thường nhưng lại vô giá đối với những ai đã từng trải qua thảm kịch.

Chỉ trong một vòng lặp trước thôi, nơi này còn chìm trong một lớp sương mù dày đặc và u ám, và thậm chí là gì đó đổ nát hoang sơ do thời gian. ``Vòng lặp thời gian'' do oán linh Naomi tạo ra đã từng là một cơn ác mộng giam cầm người dân nơi đây. Nhưng giờ đây, mọi thứ đã trở lại đúng với quỹ đạo vốn có của nó. Nhà ga Aogami đang hoạt động hết công suất. Tiếng loa phát thanh thông báo các chuyến tàu đến và đi vang lên đều đặn. Dòng người qua lại tấp nập, hối hả, trên môi nở những nụ cười rạng rỡ của một ngày mới.

Những người từng bị mất tích trong vòng lặp giờ đã trở về an toàn. Giống như những gì Kuon-san từng nói, khi nguyên nhân cốt lõi bị xóa bỏ, thế giới sẽ tự động điều chỉnh lại những lỗ hổng logic. Ký ức của toàn bộ người dân trong thị trấn đã bị viết lại một cách tinh vi. Đối với họ, những người mất tích chỉ đơn giản là vừa trở về sau một chuyến đi dài ngày, hoặc vừa xuất viện sau một cơn bạo bệnh.

Và dĩ nhiên, người dân thị trấn Aogami cũng không hề có một chút ký ức nào về sự tồn tại của ``nhà Hikawa''. Căn hộ nhỏ ở Kuon-san từng sinh sống giờ đây chỉ là một căn hộ trống vô chủ đang treo biển cho thuê. Chàng thanh niên mang mái tóc đen với ánh mắt đỏ, cô nàng tóc đuôi ngựa lúc nào cũng lăm lăm đe dọa người khác, cô nhóc hồn nhiên với nụ cười tỏa nắng, và cả linh hồn mang mái tóc tết đầy bi kịch\dots\ tất cả bọn họ đã hoàn toàn bốc hơi khỏi dòng lịch sử của thế giới này.

Nhưng lớp 1-C chúng tôi thì không quên.

Bằng một cách nào đó, có lẽ là do sự can thiệp của Game trước khi rời đi, hoặc có lẽ là do một phép màu từ thần linh ở đền Aogami, ký ức của toàn bộ học sinh lớp 1-C và cô chủ nhiệm Yuko-sensei vẫn được giữ nguyên vẹn. Chúng tôi trở thành những người duy nhất trên thế giới này mang theo những mảnh ghép của mùa hè ấy.

Hôm nay là ngày tập trung câu lạc bộ ở trường Trung học Aogami. Dù đang là kỳ nghỉ hè, khuôn viên trường vẫn khá đông đúc. Tôi đẩy nhẹ cánh cổng sắt, bước vào khoảng sân rộng lớn.

Từ phía sân điền kinh, những tiếng hô hào dõng dạc vang lên. Yae đang dẫn đầu nhóm chạy bền, mồ hôi nhễ nhại dưới cái nắng chói chang. Trong lúc đó, Inari chạy ngay phía sau, chiếc đuôi cáo giả (mà cô nàng luôn khăng khăng là hàng thật) đung đưa theo từng nhịp bước.

``Nhanh lên nào! Nếu là Kuon-san thì cậu ấy đã hoàn thành mười vòng từ đời nào rồi!'' Yae hét lớn, quay lại đốc thúc cả nhóm.

``Thế thì cậu đi mà chạy đua với cậu ấy! Mình mệt bở hơi tai rồi đây này!'' Lớp trưởng thở hồng hộc, nhưng trên môi vẫn nở một nụ cười rạng rỡ.

Họ vẫn vậy, vẫn tràn đầy năng lượng và nhiệt huyết. Nhưng cái tên ``Kuon-san'' được thốt ra một cách tự nhiên như hơi thở ấy lại khiến sống mũi tôi hơi cay cay. Sự hiện diện của anh ấy đã in sâu vào thói quen của mọi người đến mức không thể xóa nhòa.

Tôi đi dọc theo hành lang tầng một, đi ngang qua phòng Thực hành Gia chánh. Mùi thơm nức mũi của bánh bơ tỏi và thịt nướng bay ra từ khe cửa khép hờ. Xuyên qua lớp kính trong suốt, tôi nhìn thấy Saya và Wata đang bận rộn trong bếp.

Hội phó đang cẩn thận rưới lớp nước sốt đặc biệt lên những đĩa hamburger nướng. Cô bạn Thiên văn thì đang đóng gói những chiếc bánh quy vị matcha vào các túi nhỏ xinh xắn.

``Xong rồi!'' Wata vỗ tay, tự hào nhìn thành quả của mình. ``Tám phần hamburger phô mai đặc biệt. Đảm bảo năng lượng cho cả đội điền kinh sau khi tập xong!''

``Và cả bánh quy matcha nữa,'' Saya cười tươi, xếp các túi bánh lên khay. ``Còn đĩa hamburger siêu cay ngập ớt habanero này\dots\ và túi bánh quy vị dâu tây siêu ngọt này\dots''

Giọng của cô nàng chợt nhỏ dần. Cả hai cô nàng khựng lại. Bàn tay của cô lơ lửng trên không trung, trong khi ánh mắt của Wata dán chặt vào đĩa thịt nướng phủ đầy ớt đỏ au - món ăn yêu thích nhất của Kaigou, và túi bánh dâu tây - món quà vặt không thể thiếu của Suzuki.

Một sự im lặng nặng nề bao trùm lấy căn phòng bếp. Thói quen nấu ăn dư ra những khẩu phần đặc biệt ấy đã trở thành một phản xạ có điều kiện của hai người họ.

``Mình ngốc thật,'' Wata khẽ cười buồn, đưa tay lên xoa đầu. ``Họ đâu còn ở đây nữa.''

``Không sao đâu mà\~{}'' Saya vòng tay ôm lấy vai bạn, vỗ nhè nhẹ. ``Chắc chắn lúc này ở thế giới bên kia, Kaigou-san đang thèm món thịt nướng siêu cay của cậu đến phát điên lên cho xem\~{}''

Tôi lặng lẽ bước đi, không muốn phá vỡ khoảng không gian riêng tư của họ. Lớp 1-C đã chấp nhận sự thật rằng nhóm Kuon đã rời đi, nhưng nỗi nhớ thương thì vẫn cứ dai dẳng, hiển hiện trong những chi tiết nhỏ nhặt nhất của cuộc sống thường nhật.

Lên đến tầng ba, tôi bước về phía phòng phát thanh của trường. Khi vừa đến gần cửa, một giai điệu quen thuộc vang lên khiến bước chân tôi chững lại.

Đó là tiếng đệm đàn guitar acoustic mộc mạc, hòa cùng chất giọng trong trẻo, nhẹ nhàng như gió thoảng. Bài hát ``Mùa hè cuối cùng ở Aogami'' mà Yuki và Azuki đã sáng tác. Không có hệ thống loa phóng thanh ồn ào, họ chỉ đơn giản là đang ngồi hát cho nhau nghe trong phòng thu nhỏ hẹp.

Tôi khép hờ đôi mắt, tựa lưng vào mảng tường mát lạnh ngoài hành lang. Những nốt nhạc mộc mạc ấy như một sợi dây vô hình kết nối tâm hồn tôi với bầu trời xanh thẳm ngoài kia. Tôi tự hỏi, liệu ở một nơi cách xa hàng vạn năm ánh sáng, hay ở một chiều không gian không thể chạm tới\dots\ Kuon-san, Kaigou-san, Suzuki-chan và Naomi-chan có đang nhìn ngắm cùng một bầu trời với chúng tôi không? Liệu họ có đang sống những chuỗi ngày bình yên, hay lại tiếp tục lao vào những cuộc chiến khốc liệt bảo vệ trật tự vũ trụ?

Dù họ có đang ở đâu, dù tương lai có ra sao đi chăng nữa, thì sự thật là\dots\ họ đã sống trọn vẹn từng khoảnh khắc ở thị trấn Aogami này. Họ đã cười, đã khóc, đã chiến đấu, và đã yêu thương chúng tôi bằng tất cả sự chân thành.

Mùa hè rực rỡ nhất trong cuộc đời của Misora Shino này, có lẽ sẽ mãi mãi dừng lại ở độ tuổi mười lăm. Nhưng tôi không hề hối tiếc. Bởi vì tôi biết, những ký ức tuyệt đẹp ấy sẽ là hành trang tiếp thêm sức mạnh cho tôi, và cho toàn thể lớp 1-C, dũng cảm bước tiếp trên con đường trưởng thành của chính mình.

\ct

Ngày một tháng Chín. Học kỳ II của trường Trung học Aogami sau kỳ nghỉ hè đã bắt đầu.

Tôi chỉnh lại chiếc nơ trên cổ áo, hít một hơi thật sâu trước khi bước vào lớp 1-C. Bầu không khí nhộn nhịp quen thuộc lập tức bủa vây lấy tôi. Mọi người ríu rít kể cho nhau nghe về những kỳ nghỉ dưỡng cùng gia đình, những trò chơi điện tử mới phá đảo, hay những câu chuyện kỳ bí vừa sưu tầm được.

Ayame và Hanabi đang bám lấy hai cánh tay của Inari, ra sức trêu chọc cô bạn lớp trưởng cũ.

``Tuyệt vời quá đi! Hội trưởng Hội học sinh Shirayuki Inari! Nghe oai phong lẫm liệt ghê chưa!'' Hanabi cười lớn, vỗ bôm bốp vào lưng Inari.

``Tôi đã bảo nhường chức Hội trưởng cho cậu là một quyết định sáng suốt mà! Giờ thì tôi có thể thoải mái đi chơi khăm khắp trường rồi!'' Ayame cũng hùa theo, nở nụ cười tinh nghịch lộ ra hai chiếc răng khểnh.

Inari đỏ bừng cả mặt, chiếc đuôi cáo sau lưng xù lên bối rối: ``M-Mình chỉ đang cố gắng trở nên tự tin hơn thôi mà\dots\ V-Và cũng nhờ Shino-san đồng ý làm Phó hội trưởng hỗ trợ mình nữa\dots''

Tôi mỉm cười tiến lại gần, vỗ nhẹ vai cô bạn nhút nhát để giải vây cho cô bạn: ``Cậu sẽ làm được mà, Inari-san. Mình và Wata-san sẽ giúp cậu quản lý mớ rắc rối do Ayame-san gây ra.''

Tiếng chuông báo hiệu giờ vào lớp reo vang. Yuko-sensei bước vào với nụ cười điềm đạm thường ngày, tay ôm cuốn sổ điểm danh. Cả lớp nhanh chóng ổn định chỗ ngồi.

Bàn ghế trong lớp đã được sắp xếp lại cho học kỳ mới. Nhưng ở khu vực trước bàn giáo viên, vẫn có ba bộ bàn ghế trống được kê ngay ngắn cạnh nhau. Không một ai trong lớp có ý định chuyển vào ngồi đó, và Yuko-sensei cũng chưa từng yêu cầu phải dọn dẹp chúng đi.

Trên ba chiếc bàn trống ấy, Saki đã cẩn thận đặt ba lọ hoa nhỏ xinh. Một lọ hoa bỉ ngạn trắng tinh khôi, một lọ cúc vạn thọ cam rực rỡ, và một lọ hoa chuông xanh biếc. Những bông hoa được chăm sóc tỉ mỉ, mang theo sự dịu dàng của người trồng hoa, và cả nỗi nhớ thương của toàn thể lớp 1-C. Nói thật, người ngoài nhìn vào sẽ tưởng rằng ba người đó đã ra đi. Đúng là họ đã không còn ở thế giới này, nhưng không phải theo nghĩa theo cách thông thường.

Yuko-sensei mở sổ điểm danh, bắt đầu gọi tên từng người. Giọng cô vang lên đều đặn, ấm áp. Khi đến cuối danh sách, cô dừng lại một nhịp dài, ánh mắt hướng về phía ba chiếc bàn trống. Cả lớp cũng đồng loạt dõi ánh nhìn chất chứa vô vàn cảm xúc về phía góc phòng ngập tràn ánh nắng.

Nơi đó không có ai cả. Nhưng đối với lớp 1-C, dường như vẫn có một chàng trai đang lười biếng gục mặt xuống bàn ngủ gật, một cô gái đang khoanh tay liếc nhìn mọi người với ánh mắt sắc lẹm, và một cô nàng đang vừa nhai kẹo vừa tươi cười vẫy tay chào lại.

Tôi nhắm mắt lại, cảm nhận những tia nắng ấm áp của mùa thu lướt qua gò má. Mùa hè rực rỡ và kỳ diệu nhất đã chính thức khép lại, nhưng hành trình trưởng thành của chúng tôi chỉ mới vừa bắt đầu.

Ở một nơi nào đó ngoài kia, giữa vô số những giao điểm của đa vũ trụ, tôi biết chắc rằng bọn họ vẫn đang kiên cường chiến đấu. Và chúng tôi ở đây, cũng sẽ sống thật tốt phần đời của mình, để không phụ lòng những người đã từng đánh đổi tất cả vì thị trấn này.

``Kuon-san, Kaigou-san, Suzuki-san, và\dots\ Naomi-chan này\dots'' tôi thì thầm, tựa hồ muốn gửi lời chào tới nơi mà họ đang ở.

\vspace{2em}
\begin{center}
  ``Các cậu sẽ mãi là bạn của mình chứ?''
\end{center}
\vspace{2em}

\end{document}